% concatenated from FGCounter20260827f.tex
\documentclass[11pt]{article}
\usepackage{amsmath,amssymb,amsthm,mathtools}
\usepackage{geometry}
%\usepackage{hyperref}
\usepackage{enumitem}
\usepackage{microtype}
\usepackage{enumitem} % in the preamble
\usepackage{amsmath}
\usepackage[hidelinks]{hyperref}
\hypersetup{
  pdftitle={
    Counterexamples to Whole-Sequence Convergence of
    Variable-Smoothing Full-Splitting Methods
  },
  pdfauthor={Min Tao},
  pdfkeywords={
    structured fractional programming;
    variable smoothing;
    whole-sequence convergence;
    Whitney extension criterion;
    limiting lifted stationarity;
    counterexamples
  }
}

\DeclareMathOperator{\sinc}{sinc}
\geometry{margin=1in}

%------------------------------------------------
% Theorem environments
%------------------------------------------------
%\newtheorem{theorem}{Theorem}[section]
%\newtheorem{lemma}[theorem]{Lemma}
%\newtheorem{proposition}[theorem]{Proposition}
%\newtheorem{assumption}[theorem]{Assumption}
%\newtheorem{algorithm}[theorem]{Algorithm}
%\newtheorem{definition}[theorem]{Definition}
\numberwithin{equation}{section}
\newtheorem{theorem}{Theorem}[section]
\newtheorem{lemma}[theorem]{Lemma}
\newtheorem{proposition}[theorem]{Proposition}
\newtheorem{assumption}[theorem]{Assumption}

\newtheorem{remark}[theorem]{Remark}

\newtheorem{algorithm}[theorem]{Algorithm}
\newtheorem{definition}[theorem]{Definition}

\newtheorem{corollary}[theorem]{Corollary}
%------------------------------------------------
% Commands
%------------------------------------------------
\newcommand{\R}{\mathbb{R}}
\newcommand{\Hsp}{\mathcal H}
\newcommand{\ip}[2]{\left\langle #1,#2\right\rangle}
\newcommand{\norm}[1]{\left\lVert #1\right\rVert}
\newcommand{\Lip}{\operatorname{Lip}}

\newcommand{\Rop}{\mathcal R}
\newcommand{\Sop}{\mathcal S}
\newcommand{\Iop}{\mathcal I}
\newcommand{\Top}{\mathcal T}
\newcommand{\eps}{\varepsilon}
\newcommand{\Sph}{\mathbb S}
\newcommand{\Sone}{\mathbb S^1}
\newcommand{\Tone}{\mathbb T^1}
\newcommand{\clust}{\operatorname{cluster}}

\newcommand{\dist}{\operatorname{dist}}
\newcommand{\prox}{\operatorname{prox}}
\newcommand{\Proj}{\operatorname{Proj}}
\newcommand{\ind}{\iota}
\newcommand{\clB}{\overline{\mathbb B}}
\newcommand{\nn}{\nonumber}
\newcommand{\h}[1]{\mathbf{#1}}

\DeclareMathOperator{\sign}{sign}
\DeclareMathOperator{\rank}{rank}
\DeclareMathOperator{\dom}{dom}
\DeclareMathOperator*{\inte}{int}
%\hypersetup{hidelinks}
%------------------------------------------------
% Title information
%------------------------------------------------
\title{
Counterexamples to Whole-Sequence Convergence of
Variable-Smoothing Full-Splitting Methods
}

\author{
Min Tao\\[1mm]
\small School of Mathematics, Nanjing University\\
\small National Key Laboratory for Novel Software Technology\\
\small Nanjing 210093, China\\
\small \href{mailto:taom@nju.edu.cn}{taom@nju.edu.cn}
}

%\date{\today}
\date{}
\begin{document}

\maketitle
%\begin{abstract}
%We investigate whole-sequence convergence of the smoothing-based
%full-splitting proximal subgradient method (S-FSPS) for structured nonconvex
%and nonsmooth fractional programs, introduced  by
%Bo\c{t}, Li, and Tao (SIAM J. Optim., 35(4):2623--2653, 2025) as Algorithm~4.1.
%The existing convergence theory establishes only the existence of a subsequence
%converging to an exact limiting lifted stationary point, leaving open whether the
%entire generated sequence must converge. We answer this question in the negative by constructing two admissible
%instances whose primal sequences both have the cluster set
%\(\{1\}\times\mathbb S^1\) and trajectories of infinite length, although
%every cluster point is an exact limiting lifted stationary point. In the first construction, the linear operator \(A\) is nonzero and
%every generated smoothing-dual iterate \(z^k\) is nonzero. In the second, the feasible set is full-dimensional, \(A\) has full
%row rank, and each of \(f\circ K\), \(g\circ A\), and the numerator
%\(g\circ A+h\) is nonconstant on the feasible set. The first construction also
%yields a nonconvergent example for the corresponding variable-smoothing,
%single-loop, full-splitting method for nonconvex and nonsmooth composite
%optimization, even though a subsequence of the generated iterates converges to an exact  stationary
%point. Thus, a vanishing but nonsummable smoothing schedule does not ensure
%whole-sequence convergence, even when every cluster point is an exact
%stationary point.
%%To the best of our knowledge, this is the first explicit construction
%%showing that a vanishing but nonsummable smoothing-parameter sequence, even together with exact
%%stationarity, does not imply whole-sequence convergence for variable-smoothing
%%full-splitting schemes.
%\end{abstract}

\begin{abstract}
We study  whole-sequence convergence of the smoothing-based
full-splitting proximal subgradient method (S-FSPS) for structured nonconvex
and nonsmooth fractional programs, introduced  by
Bo\c{t}, Li, and Tao (SIAM J. Optim., 35(4):2623--2653, 2025) as Algorithm~4.1. Existing theory
guarantees only the existence of a subsequence converging to a
limiting lifted stationary point. We show that this guarantee is sharp by constructing two admissible
instances whose corresponding primal sequences both have cluster set
$\{1\}\times\mathbb S^1$ and infinite length, although every cluster
point is a limiting lifted stationary point.  The construction prescribes a slowly rotating spiral and
realizes it exactly through a compatible first-order jet and a
$C^{1,1}$ Whitney extension. In the first instance, $A$ has rank one
and every smoothing-dual iterate is nonzero. In the second, the
feasible set is full-dimensional, $A$ has full row rank, and each of
$f\circ K$, $g\circ A$, and the numerator $g\circ A+h$ is
nonconstant on the feasible set. The first instance also yields a nonconvergent example for the
corresponding variable-smoothing, single-loop, full-splitting method
for nonconvex and nonsmooth composite optimization, although that
method still admits a subsequence converging to an exact stationary
point. Thus, a vanishing but
nonsummable smoothing schedule does not imply whole-sequence
convergence.
\end{abstract}
{\bf Keywords:} structured fractional programming; variable smoothing; whole-sequence convergence;
 Whitney extension criterion; limiting lifted stationarity; counterexamples.
\vspace{0.2cm}
%\medskip

\textbf{Mathematics Subject Classification (2020):}
90C32, 90C26, 49M27, 65K05.




\section{Introduction}\label{sec:introduction}


Whole-sequence convergence is substantially stronger than subsequential
convergence to stationary points. In particular, boundedness of the generated
sequence and stationarity of all its cluster points do not imply convergence
to a single limit. A standard way to upgrade subsequential
convergence to whole-sequence convergence is provided by the
Kurdyka--\L ojasiewicz (KL) framework \cite{AB,ABS,Sun1}. When a suitable KL
inequality is combined with a sufficient-decrease condition, a relative-error
estimate, and appropriate control of perturbation terms, one can often
establish convergence of the entire sequence; see, for example,
\cite{AB,ABS,Sun1}. The distinction between subsequential stationarity and
whole-sequence convergence is central to the present work.


%An important objective in the convergence analysis of optimization algorithms is to
%understand not only whether stationary cluster points exist, but also whether the
%entire generated sequence converges.  The latter property is substantially stronger:
%boundedness of the iterates, convergence of the objective values, and stationarity of
%all cluster points do not ensure whole-sequence convergence.  This
%distinction is particularly relevant for nonconvex and nonsmooth methods involving a
%vanishing regularization or smoothing parameter, because the resulting approximation
%errors may vanish without being summable.

Smoothing methods for nonsmooth optimization have traditionally been analyzed
in terms of objective-value accuracy, iteration complexity, or subsequential
stationarity. Classical contributions include Nesterov's smoothing technique
\cite{Nesterov2005}, the unified framework of Beck and Teboulle
\cite{BeckTeboulle2012}, and Chen's gradient-consistent approximation framework
for nonconvex problems \cite{Chen2012}. Variable-smoothing schemes were later
developed for convex composite optimization
\cite{BotBohm2020,BotHendrich2015} and for nonconvex composite models
\cite{BohmWright2021,LiuXia2024,LongEtAl2026}. Kume and Yamada~\cite{KumeYamada2025} analyze a single-loop
variable-smoothing method for weakly convex composite minimization
with a manifold constraint via parametrization, obtaining
subsequential stationarity rather than whole-sequence convergence.
 These results provide
complexity guarantees for achieving $\varepsilon$-approximate stationarity or
accumulation-point guarantees for exact limiting (lifted) stationary points,
but do not generally establish whole-sequence convergence to an exact
stationary point.


Whole-sequence convergence is available in more specialized settings. Sabach,
Teboulle, and Voldman \cite{SabachTeboulleVoldman2018} establish finite length
and convergence for a clustering model by combining the KL property with a
fixed smoothing parameter.  Bian and Chen~\cite{BianChen2020} establish whole-sequence convergence,
without invoking the KL property, for an exact capped-$\ell_1$ relaxation
of a cardinality-penalized sparse regression model. Their proof exploits
the model-specific lower-bound and support-identification structure of the
capped-$\ell_1$ penalty, together with an adaptive smoothing update satisfying
\(\sum_{k=0}^{\infty}\mu_k^2<+\infty\). However, square summability of the smoothing parameters
alone is insufficient for whole-sequence convergence, as the counterexamples
constructed below demonstrate.

 %Their convergence analysis establishes the
%existence of a subsequence along which the corresponding smoothed-gradient
%norm converges to zero and shows that every cluster point of such a subsequence
%is an exact stationary point. These results provide
%complexity or accumulation-point guarantees, but do not generally establish
%convergence of the whole sequence to an exact stationary point.

For the structured fractional optimization problem (\ref{eq:fractional-problem}),
Bo\c{t}, Li, and Tao
\cite{BotLiTao2025} introduced the smoothing-based full-splitting proximal
subgradient method (S-FSPS).  The method separates the constituent functions and linear operators
and employs a nonsummable smoothing-parameter sequence that decreases
to zero.
The convergence theory in \cite{BotLiTao2025} guarantees, under the standing
assumptions, the existence of subsequences converging to exact limiting lifted
stationary points.  It does not establish convergence of the entire
sequence.  This leaves open whether
whole-sequence convergence of S-FSPS follows under its standing assumptions.






%The KL framework  \cite{AB,ABS,Sun1} does not directly resolve this issue. Standard KL-based convergence analyses typically derive the finite-length property by combining
% a KL property for a fixed or suitably
%augmented Lyapunov function with a sufficient-decrease property, a
%relative-error bound with uniform coefficients, and summable perturbation
%terms; see also \cite{Sun1}. These ingredients are not guaranteed for
%S-FSPS under its standing assumptions. First, these assumptions impose
%neither the KL property nor a structural condition ensuring it. Second,
%even if an appropriate KL property is assumed, the currently available smoothing-dependent estimates do not satisfy the hypotheses of the standard finite-length arguments. Moreover, even if the smoothing parameter is treated as
%an additional variable in an augmented Lyapunov function, the available
%estimates do not fit the standard KL framework. For example, suppose that
%the conjugate \(f^*\) of the function \(f\) in~\eqref{eq:fractional-problem}
%is calm \cite{RockWets} on its effective domain (as in \cite[Section 6.2]{BotLiTao2025}),
%and consider the augmented merit function
%\[
%\Upsilon: \left\{ (\h{x}, \h{y}) \in \mathbb{R}^n \times \operatorname{dom} f^* : \langle K \h{x}, \h{y} \rangle - f^*(\h{y}) > \frac{m}{2} \right\} \times (0,+\infty) \rightarrow \overline{\mathbb{R}}
%\]
%defined by
%\[
%\Upsilon(\h{x}, \h{y}, \gamma)
%= \frac{g_{\gamma}(A{\h x})+h(\h x)+\iota_{S}(\h x)}{\langle K \h{x}, \h{y} \rangle - f^*(\h{y})},
%\]
%where $m:=\inf_{x\in {\cal S}} f(Kx)>0$.
%The available relative-error estimate has the schematic form
%\[
%\begin{aligned}
%\dist\!\left(
%0,\partial\Upsilon(x^{k+1},y^{k+1},\gamma_k)
%\right)
%&\leq
%\zeta_k\|x^{k+1}-x^k\|
%+c_1\|z^{k+1}-z^k\|
%+\Delta_{k+1}
%+c_2(\gamma_{k-1}-\gamma_k),
%\end{aligned}
%\]
%where
%$
%\zeta_k\rightarrow+\infty
%\;\text{as}\;
%\gamma_k\downarrow0,
%$
%and
%$
%\Delta_{k+1}:=c_3\|z^{k+1}\|^2.
%$
%By setting
%$
%\mathcal E_{k+1}
%:=\Upsilon(x^{k+1},y^{k+1},\gamma_k),
%$
%the available descent estimate controls only
%$
%\mathcal E_{k+1}
%\leq
%\mathcal E_k
%-c_4\delta_k\|x^{k+1}-x^k\|^2
%+c_5(\gamma_{k-1}-\gamma_k),
%$
%where \(c_1,c_2,c_3,c_4,c_5>0\). Thus, \(\Delta_{k+1}\) is neither controlled by
%the descent inequality nor known to be summable. Moreover, the relative-error
%coefficient \(\zeta_k\) is nonuniform because it contains a term proportional
%to
%$
%\delta_k\asymp\gamma_k^{-1}.
%$
%Although the parameter-variation sequence
%\(\{(\gamma_{k-1}-\gamma_k)\}_{k\ge1}\) is summable, the prescribed schedule
%$
%\gamma_k\downarrow0,
%\;
%\sum_{k=0}^{\infty}\gamma_k=+\infty,
%$
%is incompatible with summability conditions such as
%\(\sum_{k=0}^{+\infty}\sqrt{\gamma_{k}}<+\infty\), which arise in the inexact KL analysis framework developed in \cite{Sun1}.
%Consequently, the available KL theory cannot be applied directly under the
%prescribed smoothing schedule to establish whole-sequence convergence of S-FSPS.
The KL framework \cite{AB,ABS,Sun1} does not directly resolve this
issue. Standard KL-based convergence analyses typically establish
the finite-length property by combining a KL inequality for a fixed
or suitably augmented Lyapunov function with a sufficient-decrease
estimate, a relative-error bound with uniform coefficients, and
summable perturbation terms. These ingredients are not guaranteed
for S-FSPS under its standing assumptions. In particular, the
assumptions impose neither the KL property nor a structural
condition, such as definability, that would ensure it.

Even if an appropriate KL property is additionally assumed, the
currently available smoothing-dependent estimates do not fit the
standard finite-length framework. As the smoothing parameter tends
to zero, the coefficients in the relative-error estimate become
nonuniform, while some residual terms are neither controlled by the
available descent inequality nor known to be summable. Treating the
smoothing parameter as an additional variable in an augmented
Lyapunov function does not remove these difficulties under the
currently available estimates. Moreover, the prescribed
nonsummable smoothing schedule is incompatible with stronger
summability conditions that arise in existing inexact KL analyses
\cite{Sun1}. Consequently, the available KL theory cannot be applied
directly to establish whole-sequence convergence of S-FSPS under its
standing assumptions.

These observations motivate us to investigate whether the primal sequence generated by S-FSPS converges. More generally, we ask whether such convergence
can be derived for variable-smoothing, single-loop, full-splitting algorithms
for broader classes of structured optimization problems.
Nonconvergence of iterates can occur even in smooth convex optimization. Bolte
and Pauwels \cite{BoltePauwels2021} give examples in which classical algorithms
and associated dynamical systems fail to converge. Bolte, Combettes, and
Pauwels \cite{BolteCombettesPauwels2024} show that Frank--Wolfe iterates may
fail to converge over a compact convex set even when their objective values
converge to the optimum. Our setting is different: the nonconvergence arises
in a variable-smoothing full-splitting method for structured fractional
optimization.

We construct a slowly rotating spiral whose angular increments tend to zero
but have a divergent sum. The orbit has the unit circle as its cluster set and
has infinite length. A compatible first-order jet and a
\(C^{1,1}\) Whitney extension are then used to construct problem data for which
S-FSPS generates this orbit exactly. Both instances satisfy the assumptions of
the subsequential convergence theory, yet their primal sequences share the cluster
set \(\{1\}\times\mathbb S^1\) and do not converge, although every cluster
point is an exact limiting lifted stationary point.
%We first refine the
%subsequential analysis of S-FSPS by showing that the explicit uniform
%subgradient bound already suffices for the relevant limiting lifted-stationarity
%conclusion, without requiring the additional interior-domain condition on
%\(A(\mathcal S)\). We then show that even this refined subsequential result
%cannot be strengthened to whole-sequence convergence under the standing
%assumptions. To this end, we construct two admissible counterexamples whose
%primal iterates trace nonconvergent slowly rotating spirals with the entire
%unit circle as their cluster set and infinite-length trajectories, although
%every cluster point is an exact limiting lifted stationary point. The first
%counterexample has a nonzero composite operator and nonzero dual iterates and,
%moreover, yields a nonconvergent example for the corresponding
%variable-smoothing, single-loop, full-splitting method for nonsmooth,
%nonconvex composite optimization. The second counterexample is a
%full-dimensional fractional instance in which \(A\) has full row rank, while both \(f\circ K\) and \(g\circ A\) are nonconstant on the
%feasible set. Thus, the failure of whole-sequence convergence persists beyond
%degenerate problem structures and is not specific to the fractional setting.

%The construction is exact rather than merely asymptotic. We prescribe the
%desired spiral together with compatible function values and first-order data
%and then invoke a \(C^{1,1}\) Whitney extension theorem to construct the smooth
%component of the objective so that the S-FSPS iteration reproduces the
%prescribed trajectory. As a further consequence, the same construction yields
%a nonconvergent example for the corresponding variable-smoothing, single-loop,
%fully splitting method for nonconvex and nonsmooth composite optimization, even
%though the generated sequence still admits a subsequence converging to an
%exact stationary point.

%A broader class of nonconvergence phenomena for first-order methods was
%investigated by Bolte and Pauwels~\cite{BoltePauwels2021}, who constructed
%counterexamples addressing several important questions in convex
%optimization. The present construction differs both in its algorithmic
%setting and in the mechanism producing nonconvergence.

\paragraph{Contributions}
The principal contributions are summarized below.
\begin{enumerate}[label=\textup{(\roman*)},leftmargin=2.3em]

\item \textbf{Sharpness of the S-FSPS convergence theory.}
We first refine the subsequential convergence analysis in
\cite[Theorem~4.3]{BotLiTao2025}. Specifically, the explicit bound
$
\operatorname{dist}\bigl(0,\partial g(Ax)\bigr)\leq \ell
\; (x\in\mathcal S)
$
already provides the uniform Moreau-envelope estimates required to establish
the existence of a subsequence converging to an exact limiting lifted
stationary point. Consequently, this subsequential conclusion remains valid
without the  condition
$
A(\mathcal S)\subseteq\operatorname{int}(\operatorname{dom}g).
$
We then show that this subsequence convergence result  cannot be strengthened to
whole-sequence convergence under the standing assumptions. To show this, we
construct an admissible instance whose primal sequence has a continuum of cluster points and an
infinite-length trajectory.


\item \textbf{Exact realization and robustness of the counterexample.}
We prescribe a slowly rotating spiral together with compatible function
values and first-order data, verify the Glaeser--Whitney compatibility
conditions globally on the union of the spiral and its limiting circle, and
invoke a \(C^{1,1}\) Whitney extension theorem
\cite{DaniilidisEtAl2018} to construct a globally defined smooth component
whose gradient has a controlled Lipschitz constant. We then
choose the remaining problem data so that the S-FSPS iteration reproduces the
prescribed orbit exactly. In the basic construction, the linear operator \(A\) is nonzero and
\(z^k\neq 0\) for every \(k\geq 0\).
 We further construct a
full-dimensional fractional instance for which \(A\) has full row rank,
while each of  \(f\circ K\), \(g\circ A\) and $g\circ A+h$ is nonconstant on the
feasible set. Thus, the failure of whole-sequence convergence cannot be
attributed to a lower-dimensional feasible set, degeneracy of the composite
operator, or constancy of the denominator \(f\circ K\), the composite term \(g\circ A\) or the numerator $g\circ A+h$.

\item \textbf{Implications for composite optimization.}
The first construction also yields a nonconvergent example for the corresponding
variable-smoothing, single-loop, full-splitting method applied to
$
\min_{x\in\mathcal S}\; g(Ax)+h(x).
$
The generated sequence still admits a subsequence converging to an exact
 stationary point, while the full sequence fails to converge.

\end{enumerate}


%\paragraph{Contributions.}
%The main contributions of this paper are as follows.
%
%\begin{enumerate}
%%\item
%%We first sharpen the subsequential convergence analysis of S-FSPS.
%%More precisely, we show that the explicit uniform subgradient bound in
%%Assumption~\ref{ass:standing} yields uniform estimates for the Moreau envelope
%%of \(g\) on \(A(\mathcal S)\). Consequently, conclusions~(vi)--(vii) of
%%\cite[Theorem~4.3]{BotLiTao2025} remain valid without the
%%assumption
%%$
%%A(\mathcal S)\subseteq\operatorname{int}(\dom g).
%%$
%
%\item
%We give a negative answer to the question of whether the standing
%assumptions of S-FSPS guarantee convergence of the entire primal sequence.
%Specifically, we construct an admissible problem for which S-FSPS generates primal iterates that follow a prescribed slowly rotating spiral satisfying
%$
%    \operatorname{cluster}(x^k)=\{1\}\times\mathbb S^1
%    \;\text{and}\;
%    \sum_{k=0}^{\infty}\|x^{k+1}-x^k\|=+\infty,
%$
%where \(\operatorname{cluster}(x^k)\) denotes the set of cluster points of
%\((x^k)_{k\geq0}\), and \(\mathbb S^1\) denotes the unit circle in
%\(\mathbb R^2\). Consequently, the primal sequence fails to converge, even though the feasible set is compact, the smoothing
%parameters vanish but are nonsummable, the composite operator \(A\) is
%nonzero, and the dual iterates remain nonzero. The prescribed orbit is
%realized exactly by constructing a compatible first-order jet and invoking
%a Whitney \(C^{1,1}\) extension theorem.
%
%
%\item
%We further show that this nonconvergence phenomenon is not confined to the
%simplified structure of the first example. A second construction can be
%chosen so that \(\mathcal S\) is full-dimensional, $A$ has full row rank and no zero rows, and both \(f\circ K\) and \(g\circ A\) are nonconstant. Hence,
%whole-sequence convergence may fail even for a fractional
%instance with nondegenerate linear operators.
%
%
%\item As a by-product, the first construction also provides a counterexample for
%the corresponding variable-smoothing full-splitting method applied to the
%nonconvex and nonsmooth composite problem
%$ \min_{x\in\mathcal S}\; g(Ax)+h(x).
%$
%The generated sequence admits a subsequence converging to an exact stationary
%point, although the entire sequence fails to converge. Therefore, the
%obstruction identified here is not specific to fractional programming but
%arises in variable-smoothing, full-splitting methods.
%\end{enumerate}
%\vspace{0.1cm}

\paragraph{Organization}
The remainder of the paper is organized as follows.
Section~\ref{sec:algorithm} recalls S-FSPS, refines its subsequential
convergence theory, and states the whole-sequence convergence question.
Section~\ref{sec:whitney} presents the Whitney extension tool used in the
construction. Section~\ref{sec:construction} develops the basic and
full-dimensional fractional counterexamples. Section~\ref{sec:com}
treats the corresponding composite optimization method. Finally,
Section~\ref{sec:conclusion} concludes the paper and discusses the implications
for whole-sequence convergence.
%\item[(i)] For the smoothing-based full-splitting method S-FSPS, we sharpen
%the existing convergence analysis. In particular, under the explicit uniform
%subgradient bound in Assumption~\ref{ass:standing}, the additional condition
%$
%A(\mathcal S)\subseteq \operatorname{int}(\operatorname{dom} g)$
%is not needed to establish conclusions~(vi)--(vii) of
%\cite[Theorem~4.3]{BotLiTao2025}.
%
%\item[(ii)] Even under the refined convergence results, whole-sequence
%convergence of S-FSPS cannot be guaranteed in general.
%
%\item[(iii)] We construct an explicit counterexample that satisfies every
%S-FSPS update exactly, while the generated primal sequence fails to converge.
%\section{Introduction}
%\section{Introduction}
%
%Whole-sequence convergence is substantially stronger than subsequential
%convergence to stationary points. In nonconvex and nonsmooth optimization, a
%generated sequence may be bounded and every one of its cluster points may be
%stationary, while the sequence itself continues to move among those cluster
%points and therefore fails to converge. One of the principal mechanisms for
%ruling out such behavior is the Kurdyka--\L ojasiewicz (KL) framework. When a
%suitable KL inequality is combined with sufficient decrease, a relative-error
%estimate, and appropriate control of perturbations, one can often establish a
%finite-length property and hence convergence of the entire sequence; see, for
%example, \cite{AB,ABS}. The gap between subsequential stationarity and
%whole-sequence convergence is central to the present work.
%
%Smoothing is a classical strategy for treating nonsmooth and nonconvex
%optimization problems. Chen~\cite{Chen2012} developed a general framework
%based on smooth approximations and gradient consistency, obtaining
%subsequential convergence to stationary points under suitable
%parameter-update rules. Building on this general smoothing philosophy,
%several algorithms have been proposed for more structured nonsmooth
%problems. One representative example is the smoothing alternating
%minimization-based algorithm (SAMBA) of Sabach, Teboulle, and
%Voldman~\cite{SabachTeboulleVoldman2018}. For a nonsmooth sum-min clustering
%model, their method introduces simplex-valued assignment variables and
%applies alternating minimization to a smoothed reformulation. Under
%boundedness and semialgebraicity assumptions, a KL analysis yields the
%finite-length property and convergence of the primal center sequence. That
%result, however, concerns a fixed smoothing parameter and therefore does not
%address whole-sequence convergence when the smoothing parameter is driven to
%zero.
%
%The variable-smoothing framework of Kume and Yamada~\cite{KumeYamada2025}
%illustrates the additional difficulty created by a vanishing smoothing
%parameter. Under their standing assumptions, the sequence generated by
%Algorithm~1 possesses a subsequence converging to an exact stationary point,
%but convergence of the entire sequence is not established. Thus, in
%variable-smoothing methods, the passage from subsequential stationarity to
%whole-sequence convergence remains a genuinely nontrivial issue.
%
%Fractional optimization has likewise motivated an extensive literature on
%first-order and proximal algorithms. Bo\c{t}, Dao, and
%Li~\cite{BotDaoLi2022} developed extrapolated proximal subgradient methods for
%a class of nonsmooth and nonconvex fractional programs, while Zhang and
%Li~\cite{ZhangLi2022} proposed proximity-gradient-subgradient algorithms and
%line-search variants for another class of fractional optimization problems.
%More recently, Bo\c{t}, Li, and Tao~\cite{BotLiTao2025} studied the structured
%fractional model~\eqref{eq:fractional-problem}, in which nonsmooth terms in
%the numerator and denominator are composed with linear operators. Since the
%proximal mappings of the resulting composite functions are generally
%difficult to evaluate, they proposed a family of single-loop, fully splitting
%proximal subgradient methods requiring only the proximal mappings of the
%component functions or their Fenchel conjugates.
%Of particular relevance here is the smoothing-based fully splitting proximal
%subgradient method, abbreviated as S-FSPS, which employs a vanishing
%smoothing parameter while retaining the fully splitting structure. The
%convergence theory in \cite[Theorem~4.3]{BotLiTao2025} guarantees the
%existence of subsequences converging to exact limiting lifted stationary
%points. It does not establish convergence of the entire primal sequence. This
%leaves the following natural question: do the standing assumptions of S-FSPS
%nevertheless guarantee whole-sequence convergence?
%
%We answer this question in the negative. We construct an admissible problem
%instance for which the primal iterates generated by S-FSPS exactly follow a
%slowly rotating spiral whose cluster set is a full unit circle. More
%precisely, in the full-dimensional construction, the cluster set is
%\(\{1\}\times\mathbb S^1\), where \(\mathbb S^1\) denotes the unit circle in
%\(\mathbb R^2\). Consequently, the primal sequence does not converge and its
%trajectory has infinite length. Nevertheless, every cluster point is an
%exact limiting lifted stationary point, and all the standing assumptions used
%in the existing subsequential convergence theory are satisfied. The
%subsequential convergence result of \cite[Theorem~4.3]{BotLiTao2025} is
%therefore essentially sharp under those assumptions.
%
%The nonconvergence phenomenon is not caused by the most immediate
%degeneracies of the model. A strengthened construction uses a
%full-dimensional feasible set and a full-row-rank composite operator with no
%zero rows, while both the numerator and the denominator of the fractional
%objective are nonconstant. Thus, neither vanishing smoothing parameters nor
%exact limiting lifted stationarity of all cluster points is sufficient to
%ensure convergence of the whole sequence, even in a structurally nondegenerate
%instance.
%
%The construction is exact rather than asymptotic. We first prescribe the
%desired spiral trajectory together with compatible function values and
%first-order data. We then invoke a \(C^{1,1}\) Whitney extension theorem to
%construct the smooth component of the objective so that each S-FSPS update
%returns exactly the next point of the prescribed trajectory. This realization
%of a prescribed nonconvergent orbit is the principal technical ingredient of
%the paper. As a further consequence, the same construction yields a
%nonconvergent instance for the corresponding variable-smoothing, single-loop,
%fully splitting method for nonconvex and nonsmooth composite optimization. In
%that instance, the generated sequence again possesses subsequences converging
%to exact limiting stationary points but does not converge as a whole.
%
%The negative result is consistent with the standard KL convergence theory.
%Indeed, the standing assumptions of S-FSPS do not impose a KL or definability
%condition on the problem data. More importantly, the difficulty is not merely
%that the merit function depends on the smoothing parameter
%\(\gamma_k\). One could, in principle, attempt to treat \(\gamma\) as an
%additional variable and construct an augmented Lyapunov function. The
%estimates currently available for S-FSPS, however, do not directly satisfy the
%hypotheses of the standard KL frameworks
%\cite{AB,ABS,Sun1}. Specifically, the available relative-error estimate
%contains an additional term that is not controlled by the corresponding
%descent inequality, while a coefficient multiplying the step length becomes
%unbounded as \(\gamma_k\downarrow0\). Moreover, the prescribed smoothing
%schedule
%\[
%    \gamma_k\downarrow0,
%    \qquad
%    \sum_{k=0}^{\infty}\gamma_k=+\infty,
%\]
%does not by itself imply the summability of the perturbation terms required by
%existing inexact KL analyses such as \cite{Sun1}. Consequently, those
%frameworks cannot be invoked to deduce a finite-length property under the
%standing assumptions of S-FSPS. Our counterexample shows that this failure is
%not merely a limitation of the currently available proof techniques: without
%additional assumptions, whole-sequence convergence is in fact false.
%
%A broader family of nonconvergence phenomena for first-order methods was
%investigated by Bolte and Pauwels~\cite{BoltePauwels2021}, who constructed
%counterexamples addressing several fundamental questions in convex
%optimization. The present work is complementary to theirs. It concerns a
%variable-smoothing, fully splitting method for structured fractional and
%composite optimization, and the mechanism generating nonconvergence is
%different: the prescribed spiral is reproduced exactly by the algorithmic
%iteration, rather than arising from an asymptotic approximation or from
%surrogate dynamics.
%Whole-sequence convergence is substantially stronger than subsequential
%convergence to stationary points. In particular, boundedness of the generated
%sequence and stationarity of all its cluster points do not, by themselves,
%imply convergence to a single limit. A standard mechanism for upgrading
%subsequential convergence to whole-sequence convergence is provided by the
%Kurdyka--\L ojasiewicz (KL) framework. When a suitable KL inequality is
%combined with a sufficient-decrease condition, a relative-error estimate, and
%appropriate control of perturbation terms, one can often establish the
% convergence of
%the entire sequence; see, for example,
%\cite{AB,ABS}. The distinction between
%subsequential stationarity and whole-sequence convergence is central to the
%present work.
%
%Smoothing is a classical strategy for addressing nonsmooth and nonconvex
%optimization problems. Chen~\cite{Chen2012} developed a general framework
%based on smooth approximations and gradient consistency, establishing
%subsequential convergence to stationary points under suitable parameter-update
%rules. Building on this general smoothing philosophy, several algorithmic
%frameworks have been developed for more structured nonsmooth problems. An
%instructive example is the smoothing alternating minimization-based algorithm
%(SAMBA) of Sabach, Teboulle, and
%Voldman~\cite{SabachTeboulleVoldman2018}. Their method reformulates a
%nonsmooth sum-min clustering problem by introducing simplex-valued assignment
%variables and then applies alternating minimization to the resulting smoothed
%problem. Under boundedness and semialgebraicity assumptions, a KL-based
%analysis establishes the finite-length property and convergence of the primal
%center sequence. This result, however, is obtained for a fixed smoothing
%parameter and therefore does not address whole-sequence convergence when the
%smoothing parameter is driven to zero.
%
%A related issue arises in the variable-smoothing framework of Kume and
%Yamada~\cite{KumeYamada2025}. Under the corresponding standing assumptions,
%the sequence generated by their Algorithm~1 is shown to possess a subsequence
%converging to an exact stationary point, whereas convergence of the entire
%sequence is not established. Thus, for variable-smoothing methods, the passage
%from subsequential stationarity to whole-sequence convergence remains a
%nontrivial issue.
%
%Fractional optimization has likewise motivated a substantial literature on
%first-order and proximal algorithms. Bo\c{t}, Dao, and
%Li~\cite{BotDaoLi2022} developed extrapolated proximal subgradient methods for
%a class of nonsmooth and nonconvex fractional programs, while Zhang and
%Li~\cite{ZhangLi2022} proposed proximity-gradient-subgradient algorithms and
%line-search variants for a class of fractional optimization problems. More
%recently, Bo\c{t}, Li, and Tao~\cite{BotLiTao2025} considered the structured
%fractional model~\eqref{eq:fractional-problem}, in which nonsmooth terms in the
%numerator and denominator are composed with linear operators. This composite
%structure makes direct proximal treatment difficult and naturally motivates
%a fully splitting approach.
%
%Accordingly, Bo\c{t}, Li, and Tao~\cite{BotLiTao2025} proposed a family of
%single-loop, fully splitting proximal subgradient methods that require only
%evaluations of the proximal mappings of the component functions or their
%conjugates, rather than those of the corresponding composite mappings. In
%particular, they introduced a smoothing-based variant, referred to as S-FSPS,
%which employs a vanishing smoothing parameter while preserving the fully
%splitting structure. The convergence analysis in
%\cite[Theorem~4.3]{BotLiTao2025} guarantees the existence of subsequences
%converging to exact limiting lifted stationary points. It does not, however,
%establish convergence of the entire primal sequence. This leaves open whether
%whole-sequence convergence of S-FSPS follows from its standing assumptions.
%
%Standard KL convergence arguments
%\cite{AB,ABS} typically combine a KL
%inequality for a fixed, or suitably augmented, Lyapunov function with a
%sufficient-decrease condition, a relative-error estimate, and summable
%perturbation terms. Under the standing assumptions of S-FSPS, the available
%convergence analysis does not provide this complete collection of estimates
%uniformly as the smoothing parameter tends to zero. In particular, the
%vanishing smoothing parameter introduces nonuniform terms that are not
%controlled by the available descent estimate.
%
%For S-FSPS, the obstruction is more subtle than the mere dependence of the
%merit function on the smoothing parameter \(\gamma_k\). One could, in
%principle, attempt to incorporate \(\gamma\) as an additional variable in an
%augmented Lyapunov function. The main difficulty is that the estimates
%currently available for S-FSPS do not directly satisfy the hypotheses of
%standard KL frameworks
%\cite{AB,ABS,Sun1}. More precisely,
%the relative-error estimate contains an additional term that has no
%counterpart in the corresponding descent inequality, while the coefficient
%multiplying the step length becomes unbounded as
%\(\gamma_k\downarrow0\). Moreover, the prescribed smoothing schedule
%\[
%    \gamma_k\downarrow0,
%    \qquad
%    \sum_{k=0}^{\infty}\gamma_k=+\infty
%\]
%does not, by itself, provide the summability conditions required by existing
%inexact KL analyses such as \cite{Sun1}. Consequently, the currently
%available KL frameworks do not yield a finite-length property and therefore
%do not establish whole-sequence convergence of S-FSPS under its standing
%assumptions.
%
%In this note, we show that the answer to the whole-sequence convergence
%question is negative. We construct an admissible problem instance for which
%the primal iterates generated by S-FSPS exactly trace a slowly rotating spiral
%whose cluster set is the entire unit circle. The primal sequence therefore
%fails to converge, and its trajectory has infinite length, despite satisfying
%all the standing assumptions underlying the existing subsequential convergence
%theory. Moreover, every cluster point is an exact limiting lifted stationary
%point. Thus, neither the vanishing of the smoothing parameters nor exact
%limiting lifted stationarity of all cluster points is sufficient to guarantee
%whole-sequence convergence.
%
%The construction is exact rather than merely asymptotic. We prescribe the
%desired spiral together with compatible function values and first-order data,
%and then invoke a \(C^{1,1}\) Whitney extension theorem to construct the smooth
%component of the objective so that the S-FSPS iteration reproduces the
%prescribed trajectory exactly. As a further consequence, the same construction
%yields a nonconvergent example for the corresponding variable-smoothing,
%single-loop, fully splitting method for nonsmooth nonconvex composite
%optimization, even though the generated sequence still possesses a subsequence
%converging to an exact limiting stationary point.
%
%A broader class of nonconvergence phenomena for first-order methods was
%investigated by Bolte and Pauwels~\cite{BoltePauwels2021}, who constructed
%counterexamples addressing several fundamental questions in convex
%optimization. The present construction differs both in its algorithmic
%setting and in the mechanism producing nonconvergence. In particular, the
%spiral trajectory does not arise from an asymptotic approximation or from
%surrogate dynamics; instead, it is generated exactly by the prescribed
%S-FSPS iteration.

%A broader class of nonconvergence phenomena for first-order methods has been
%investigated by Bolte and Pauwels~\cite{BoltePauwels2021}, who constructed
%counterexamples addressing several important questions in convex
%optimization. Our construction differs both in the algorithmic framework and
%in the mechanism responsible for nonconvergence. In particular, the spiral
%orbit is not obtained through an asymptotic approximation or surrogate
%dynamics; rather, it is generated exactly by the prescribed S-FSPS updates.
%
%\vspace{0.2cm}

%\paragraph{Contributions.}
%The main contributions of this paper are as follows.
%
%\begin{enumerate}
%%\item
%%We first sharpen the subsequential convergence analysis of S-FSPS.
%%More precisely, we show that the explicit uniform subgradient bound in
%%Assumption~\ref{ass:standing} yields uniform estimates for the Moreau envelope
%%of \(g\) on \(A(\mathcal S)\). Consequently, conclusions~(vi)--(vii) of
%%\cite[Theorem~4.3]{BotLiTao2025} remain valid without the
%%assumption
%%$
%%A(\mathcal S)\subseteq\operatorname{int}(\dom g).
%%$
%
%\item
%We give a negative answer to the question of whether the standing
%assumptions of S-FSPS guarantee convergence of the entire primal sequence.
%Specifically, we construct an admissible problem for which the S-FSPS
%iterations generate a prescribed slowly rotating spiral satisfying
%$
%    \operatorname{cluster}(x^k)=\{1\}\times\mathbb S^1
%    \;\text{and}\;
%    \sum_{k=0}^{\infty}\|x^{k+1}-x^k\|=+\infty,
%$
%where \(\operatorname{cluster}(x^k)\) denotes the set of cluster points of
%\((x^k)_{k\geq0}\), and \(\mathbb S^1\) denotes the unit circle in
%\(\mathbb R^2\). Consequently, the primal sequence fails to converge and
%has infinite length, even though the feasible set is compact, the smoothing
%parameters vanish but are nonsummable, the composite operator \(A\) is
%nonzero, and the dual iterates remain nonzero. The prescribed orbit is
%realized exactly by constructing a compatible first-order jet and invoking
%a Whitney \(C^{1,1}\) extension theorem.
%
%
%\item
%We further show that this nonconvergence phenomenon is not confined to the
%simplified structure of the first example. A second construction can be
%chosen so that \(\mathcal S\) is full-dimensional, $A$ has full row rank and no zero rows, and both \(f\circ K\) and \(g\circ A\) are nonconstant. Hence,
%whole-sequence convergence may fail even for a fractional
%instance with nondegenerate linear operators.
%
%
%\item As a by-product, the first construction also provides a counterexample for
%the corresponding smoothing-based fully splitting method applied to the
%nonconvex and nonsmooth composite problem
%\[
%    \min_{x\in\mathcal S}\; g(Ax)+h(x).
%\]
%The generated sequence admits a subsequence converging to an exact stationary
%point, although the entire sequence fails to converge. Therefore, the
%obstruction identified here is not specific to fractional programming but
%arises more broadly in variable-smoothing, single-loop splitting schemes.
%\end{enumerate}
%\vspace{0.1cm}
%%\item[(i)] For the smoothing-based full-splitting method S-FSPS, we sharpen
%%the existing convergence analysis. In particular, under the explicit uniform
%%subgradient bound in Assumption~\ref{ass:standing}, the additional condition
%%$
%%A(\mathcal S)\subseteq \operatorname{int}(\operatorname{dom} g)$
%%is not needed to establish conclusions~(vi)--(vii) of
%%\cite[Theorem~4.3]{BotLiTao2025}.
%%
%%\item[(ii)] Even under the refined convergence results, whole-sequence
%%convergence of S-FSPS cannot be guaranteed in general.
%%
%%\item[(iii)] We construct an explicit counterexample that satisfies every
%%S-FSPS update exactly, while the generated primal sequence fails to converge.





%\vspace{0.1cm}
%\noindent\textbf{Organization.}
%The remainder of the paper is organized as follows.
%Section~\ref{sec:algorithm} recalls S-FSPS, refines the subsequential
%convergence results in \cite[Theorem~4.3]{BotLiTao2025}, and states the
%open question addressed in this paper.
%Section~\ref{sec:whitney} presents the Whitney extension tool used in
%our construction.
%Section~\ref{sec:construction} develops the basic and full-dimensional
%fractional counterexamples.
%Section~\ref{sec:com} establishes the corresponding subsequential
%results for composite nonconvex optimization and specializes the
%counterexample to that setting.
%Finally, Section~\ref{sec:conclusion} concludes the paper.
%The remainder of the paper is organized as follows.
%Section~\ref{sec:algorithm}  recalls S-FSPS and refines its subsequential convergence theory in \cite[Theorem~4.3]{BotLiTao2025}, and formulates the open question addressed in this paper. Section~\ref{sec:whitney} presents the Whitney extension tool.  Section~\ref{sec:construction} constructs the basic and full-dimensional fractional counterexamples. Section \ref{sec:conclusion} derives the corresponding composite-optimization subsequential convergence results and counterexample, and Section~\ref{sec:conclusion}  concludes the paper.
%Section~\ref{sec:algorithm} recalls the S-FSPS iteration and its standing assumptions, refines the convergence results in \cite[Theorem~4.3]{BotLiTao2025}, and formulates the open question addressed in this paper. Section~\ref{sec:whitney} presents the Whitney extension theorem and a practical variant used in our construction. Section~\ref{sec:construction} develops first a basic counterexample and then a full-dimensional fractional one. Finally, Section~\ref{sec:conclusion} summarizes the implications and limitations of our results.



\section{S-FSPS and the Whole-Sequence Convergence Question}
\label{sec:algorithm}
 The Euclidean norm is denoted by \(\|\cdot\|\), and
\(\langle\cdot,\cdot\rangle\) denotes the corresponding inner product.
Let $\overline \R:=\R\cup\{+\infty\}$.
For a function \(f:\mathbb{R}^n\to\overline{\mathbb{R}}\), its
\emph{effective domain} is defined by
$
\dom f:=\{x\in\mathbb{R}^n:f(x)<+\infty\},
$
and \(f\) is said to be \emph{proper} if \(\dom f\neq\emptyset\).
The \emph{(Fenchel) conjugate} of \(f\) is given by
$
f^*(v):=\sup_{x\in\mathbb{R}^n}
\{\langle v,x\rangle-f(x)\}.
$
The domain of the convex subdifferential is denoted by
$
\dom(\partial f):=\{x\in\mathbb{R}^n:\partial f(x)\neq\emptyset\}.
$
For a proper, convex, and lower semicontinuous function
\(f:\mathbb{R}^n\to\overline{\mathbb{R}}\), its \emph{proximal operator}
with parameter \(\gamma>0\) is defined by
$
\prox_{\gamma f}(x)
:=
\arg\min_{y\in\mathbb{R}^n}
\left\{
f(y)+\frac{1}{2\gamma}\|y-x\|^2
\right\}.
$
The corresponding \emph{Moreau envelope} is defined by
$
f_\gamma(x)
:=
\min_{y\in\mathbb{R}^n}
\left\{
f(y)+\frac{1}{2\gamma}\|y-x\|^2
\right\}.
$
For every \(x\in\mathbb{R}^n\),
$
f_\gamma(x)
=
\left(f^*+\frac{\gamma}{2}\|\cdot\|^2\right)^*(x),
$
and \(f_\gamma\) is continuously differentiable with
$
\nabla f_\gamma(x)
=
\frac{x-\prox_{\gamma f}(x)}{\gamma}
=
\prox_{f^*/\gamma}\left(\frac{x}{\gamma}\right).
$
Moreover,
$
\nabla f_\gamma(x)
\in
\partial f\bigl(\prox_{\gamma f}(x)\bigr).
$
We use
$
    \Sone:=\bigl\{x\in\mathbb R^2:\|x\|=1\bigr\}
    \;\text{and}\;
    \overline{\mathbb B}_{\mathbb R^2}
    :=\bigl\{x\in\mathbb R^2:\|x\|\leq1\bigr\}
$
to denote the unit circle and the closed unit ball in \(\mathbb R^2\),
respectively. When the  Euclidean space is clear from the context,
we  write \(\overline{\mathbb B}\) for the corresponding closed unit
ball centered at the origin.
For a nonempty closed convex set
\(\mathcal C\subseteq\mathbb{R}^n\), its \emph{indicator function} is defined by
\[
\iota_{\mathcal C}(x)
:=
\begin{cases}
0, & \text{if }x\in\mathcal C,\\
+\infty, & \text{otherwise}.
\end{cases}
\]
The \emph{normal cone} to \(\mathcal C\) at \(x\) is denoted by
$
N_{\mathcal C}(x):=\partial\iota_{\mathcal C}(x)
$, and
$\Proj_{\mathcal C}(\cdot)$ denotes the Euclidean projection onto $\mathcal C$.
For a sequence \(\{x^k\}_{k\ge0}\), we denote by
$
\operatorname{cluster}(x^k)
$
the set of its cluster points.
Given a linear operator $A: {\mathbb R}^n \rightarrow {\mathbb R}^m$, we denote by $A^*: {\mathbb R}^m \rightarrow {\mathbb R}^n$ its {\it adjoint operator}. We also use $\sigma_A:=\|A\|=\sup\{ \|A  x\|: \| x\|=1\}$ to denote its norm.
Given a vector \(x\in\mathbb{R}^n\), \(x_i\) denotes its \(i\)th component.

%\subsection{S-FSPS and the Whole-Sequence Convergence Question}

In this paper, we consider the structured fractional optimization problem
\begin{equation}\label{eq:fractional-problem}
    \min_{x\in\mathcal S}
    \frac{g(Ax)+h(x)}{f(Kx)},
\end{equation}
where \(\mathcal S\) is a nonempty, convex, and compact subset of
\(\mathbb R^n\);
\(f:\mathbb R^p\to\overline{\mathbb R}:=\mathbb R\cup\{+\infty\}\) and
\(g:\mathbb R^s\to\overline{\mathbb R}\) are proper, convex, and lower
semicontinuous;
\(A:\mathbb R^n\to\mathbb R^s\) and
\(K:\mathbb R^n\to\mathbb R^p\) are linear operators; and
\(h:\mathbb R^n\to\mathbb R\) is  differentiable on an open set containing ${\cal S}$  and its gradient $\nabla h$ is $L_{\nabla h}$-Lipschitz continuous over this  set.


%In this paper, we consider the structured fractional optimization problem
%\begin{equation}\label{eq:fractional-problem}
%    \min_{x\in\mathcal S}
%    \frac{g(Ax)+h(x)}{f(Kx)},
%\end{equation}
%where \(\mathcal S\) is a nonempty, convex, and compact subset of
%\(\mathbb R^n\);
%\(f:\mathbb R^p\to\overline{\mathbb R}:=\mathbb R\cup\{+\infty\}\) and
%\(g:\mathbb R^s\to\overline{\mathbb R}\) are proper, convex, and lower
%semicontinuous;
%\(A:\mathbb R^n\to\mathbb R^s\) and
%\(K:\mathbb R^n\to\mathbb R^p\) are linear operators; and
%\(h:\mathbb R^n\to\mathbb R\) is  differentiable on an open set containing ${\cal S}$  and its gradient $\nabla h$ is $L_{\nabla h}$-Lipschitz continuous over this  set.


We use the following  assumption of \cite[Assumption~3.1]{BotLiTao2025}.
\begin{assumption}\label{ass:standing}
The set \(\mathcal S\subseteq\mathbb R^n\) is nonempty, convex, and compact;
\(A\) and \(K\) are nonzero linear operators; \(g\) and \(f\) are proper,
convex, and lower semicontinuous functions; and \(h\) is differentiable with an
\(L_{\nabla h}\)-Lipschitz continuous gradient on an open neighborhood of
\(\mathcal S\). Moreover,
\begin{enumerate}[label=\textup{(\roman*)}]
    \item \(K(\mathcal S)\subseteq\operatorname{int}(\dom f)\) and
          \(f(Kx)>0\) for every \(x\in\mathcal S\);
    \item \(\mathcal S\cap A^{-1}(\dom g)\neq\varnothing\) and
          \(\inf_{x\in\mathcal S}\{g(Ax)+h(x)\}>0\);
    \item \(A(\mathcal S)\subseteq\dom(\partial g)\), and there exists
          \(\ell>0\) such that
          \(\dist(0,\partial g(Ax))\le\ell\) for every
          \(x\in\mathcal S\).
\end{enumerate}
\end{assumption}
%\begin{remark} \(\mathcal S\cap A^{-1}(\dom g)\neq\varnothing\) is redundant since \(A(\mathcal S)\subseteq\dom(\partial g)\subseteq\dom g\) is imposed. If $g$ is globally Lipschitz continuous, we can remove Assumption \ref{ass:standing}(iii). \end{remark}
\begin{remark}
The condition
$
\mathcal S\cap A^{-1}(\dom g)\!\neq\!\varnothing
$ in Assumption~\ref{ass:standing}(ii)
is redundant under Assumption~\ref{ass:standing}(iii). Since
$
A(\mathcal S)\subseteq\dom(\partial g)\subseteq\dom g,
$
we have
$
\mathcal S\subseteq A^{-1}(\dom g),
$
which immediately yields
$
\mathcal S\cap A^{-1}(\dom g)=\mathcal S\neq\varnothing.
$
Moreover, if \(g:\R^s\rightarrow\R\) is globally Lipschitz continuous,
Assumption~\ref{ass:standing}(iii) is automatically satisfied. Indeed, if
\(L:=\Lip(g)\), then
$\partial g(u)\subseteq L\,\overline{\mathbb B}$ for all $u\in\R^s$.
\end{remark}
We review the concept of limiting  lifted stationary point of (\ref{eq:fractional-problem}).

\begin{definition}\label{def:limiting-lifted-stationarity}
For problem~\eqref{eq:fractional-problem}, a point
\[
x\in\mathcal S\cap A^{-1}(\dom\partial g)
\cap K^{-1}(\dom\partial f)
\]
with \(f(Kx)>0\) is called a \emph{limiting lifted stationary point} if
\[
0\in
f(Kx)
\bigl(
A^*\partial g(Ax)+\nabla h(x)+\partial\iota_{\mathcal S}(x)
\bigr)
-
(g(Ax)+h(x))K^*\partial f(Kx).
\]
\end{definition}

%\begin{remark}
%Assumption~\ref{ass:standing} coincides with
%\cite[Assumption~3.1]{BotLiTao2025}, except that we omit the second part of
%condition~(f), which requires the existence of a constant $\ell>0$ satisfying
%\[
%\operatorname{dist}(0,\partial g(Ax))\le \ell,
%\qquad \forall x\in\mathcal S .
%\]
%This boundedness condition is not required for the well-definedness of
%S-FSPS (\cite[Algorithm~4.1]{BotLiTao2025}) or for its convergence analysis
%in \cite[Theorem~4.3]{BotLiTao2025}. It is only used in the analysis of the
%adaptive variant proposed in \cite[Algorithm~5.1]{BotLiTao2025}. Therefore,
%we exclude this additional requirement from our standing assumptions.
%\end{remark}

\noindent Next, we define the smoothed potential function with $\gamma>0$,
\begin{equation}\label{eq:Psi}
\Psi(x,z;\gamma)
:=
\langle z,Ax\rangle-g^*(z)+h(x)+\iota_{\mathcal S}(x)
-\frac{\gamma}{2}\|z\|^2 .
\end{equation}

%\noindent For $\gamma>0$, denote the Moreau envelope and proximal mapping of $g$ by
%\[
%g_\gamma(w)
%:=\inf_{u\in\R^s}
%\left\{g(u)+\frac{1}{2\gamma}\|u-w\|^2\right\},
%\qquad
%\prox_{\gamma}(w):=\prox_{\gamma g}(w).
%\]



The following lemma shows that the explicit uniform subgradient bound in Assumption~\ref{ass:standing}(iii) provides all the local control needed in the proof of Theorem~\ref{PriTheo2R}.
\begin{lemma}
\label{lem:uniform-moreau-estimates}
Suppose that $g$ is proper, convex, and lower semicontinuous and that Assumption \ref{ass:standing}(iii) holds.
Then, for every $w\in A(\mathcal S)$ and every $\gamma>0$,
\begin{itemize}
\item[(i)]
$\|\nabla g_\gamma(w)\|\leq\ell$;
\label{eq:uniform-gradient-bound}
\item[(ii)]
$
\|\prox_{\gamma g}(w)-w\|\leq\gamma\ell$;
\label{eq:uniform-proximal-bound}
\item[(iii)]
$0\leq g(w)-g_\gamma(w)\leq\frac{\gamma\ell^2}{2}$.

\item[(iv)]
$|g(u)-g(v)|\leq\ell\|u-v\|
\;\;\text{for any }u,v\in A(\mathcal S).$
\end{itemize}
\end{lemma}
The proof of Lemma \ref{lem:uniform-moreau-estimates} is deferred to Appendix \ref{appA}.


\noindent We next recall the S-FSPS algorithm proposed in  \cite[Algorithm~4.1]{BotLiTao2025}.
\begin{algorithm}[S-FSPS]\label{alg:sfsps}
Let $\{\gamma_k\}_{k\ge0}$ be positive and nonincreasing, with
\begin{eqnarray}\label{eq:gammak}
\gamma_k\to0,
\qquad
\sum_{k=0}^{\infty}\gamma_k=+\infty.
\end{eqnarray}
Choose $\chi>1$, $\theta_0>0$, and a positive upper bound $L_{\nabla h}$ on the Lipschitz constant of
$\nabla h$ over the prescribed neighborhood, and
\begin{equation}\label{eq:delta-general}
\delta_k:=\chi\left(L_{\nabla h}+\frac{\sigma_A^2}{\gamma_k}\right).
\end{equation}
Starting from $(x^0,z^0)$ with $x^0\in\mathcal S$, generate, for $k\ge0$,
\begin{align}
 y^{k+1}&\in\partial f(Kx^k),\label{eq:y-update}\\
 x^{k+1}&:=\Proj_{\mathcal S}\left(
 x^k+\frac{\theta_k}{\delta_k}K^*y^{k+1}
 -\frac{1}{\delta_k}\nabla h(x^k)
 -\frac{1}{\delta_k}A^*z^k\right),\label{eq:x-update}\\
 z^{k+1}&:=\prox_{g^*/\gamma_k}\left(\frac{Ax^{k+1}}{\gamma_k}\right),\label{eq:z-update}\\
 \theta_{k+1}&:=\frac{\Psi(x^{k+1},z^{k+1};\gamma_k)}{f(Kx^{k+1})}.
 \label{eq:theta-update}
\end{align}
\end{algorithm}
We next refine the convergence analysis of
Algorithm~\ref{alg:sfsps} presented in \cite{BotLiTao2025}.
Specifically, we show that conclusions~(vi)--(vii) of
\cite[Theorem~4.3]{BotLiTao2025} remain valid without the assumption
$
    A(\mathcal{S})\subseteq\operatorname{int}(\dom g).$
%For convenience, define
%$
%    $
\begin{theorem}
\label{PriTheo2R}
Suppose Assumption \ref{ass:standing} holds. Define $V^k := (x^k,y^k,z^k), \; k\geq 1$. Let $\Omega$ be the set of  cluster points of the sequence $\{{ V}^k\}$ generated by Algorithm \ref{alg:sfsps}.  Then the following statements hold:
\begin{itemize}
\item[(i)] For every $k \geq 1$, the following inequality holds:
\begin{align*}\label{desineqRR}
& \ \Psi({ x}^{k+1}, { z}^{k+1}; \gamma_k)
+ \theta_k \left[f(K{ x}^k)-\left( \langle K{ x}^{k+1},{ y}^{k+1} \rangle - f^*({ y}^{k+1})\right)\right]\nn \\
\leq & \ \Psi({ x}^{k}, { z}^{k}; \gamma_{k-1})
-{\widetilde c}_{k}\|{ x}^k-{ x}^{k+1}\|^2 + \Xi^{k+1},
\end{align*}
 where
$
\Xi^{k+1} := \frac{\gamma_{k-1}-\gamma_{k}}{2}\|{ z}^{k+1}\|^2 \geq 0 \quad \mbox{and} \quad
 {\widetilde c}_{k}:=\frac{(\chi-1)}{2}\left(L_{\nabla h} +\frac{\sigma_A^2}{\gamma_{k}}\right) >0.
$

\item[(ii)] The sequence $\{{ V}^k\}$ is bounded.
\item[(iii)] There exists an index $K_1 \geq 1$ such that $\theta_k\ge 0$ for all $k\ge K_1$.
\item[(iv)] $\lim_{k\to+\infty}\theta_k=\overline{\theta}$ {for some $\overline{\theta} \ge 0$.}
\item[(v)] We have $\liminf\limits_{k\to+\infty}\delta_{k}\|{ x}^{k+1}-{ x}^k\|=0.$
\end{itemize}
%\ \ \ \   {Below, we further assume that $A({\cal S}) \subseteq \text{\rm int}(\dom  g)$.}
\begin{itemize}
\item[(vi)]
{For every $(\overline{ x},{\overline{ y}},{\overline{ z}}) \in \Omega$, we have }
$\frac{g(A{\overline{ x}})+h({\overline{ x}})+\iota_{\cal S}(\overline{ x})}{f(K{\overline{ x}})}={\overline{ \theta}},$
{where $\overline{\theta}$ is given as in {(iv)}.}
\item[(vii)]
Let \(\{{ x}^{k_j}\}\) be a subsequence of $\{x^k\}_{k\ge 0}$ such that
$
\lim\limits_{j \to +\infty} \delta_{k_j} \|{ x}^{k_j+1} - { x}^{k_j}\| = 0$ (whose existence is guaranteed by {(v)}).  Then every cluster point \(\overline{ x}\) of this subsequence is a limiting lifted stationary point for the optimization problem \eqref{eq:fractional-problem}.
\end{itemize}
\end{theorem}
The proof of Theorem~\ref{PriTheo2R} is deferred to Appendix~\ref{appB}.

\medskip
\noindent
Theorem~\ref{PriTheo2R} guarantees the existence of a subsequence converging to an exact limiting lifted stationary point. However, it remains {\it open} whether the entire primal sequence generated by S-FSPS necessarily converges.
Section \ref{sec:construction} answers this question in the negative.


\section{Whitney Extension Criterion}\label{sec:whitney}
We first recall a Glaeser-Whitney \(C^{1,1}\) extension theorem \cite{DaniilidisEtAl2018} and then derive a practical extension criterion.
Let $S$ be a nonempty subset of a Hilbert space
$(\mathcal H,\langle\cdot,\cdot\rangle,\|\cdot\|)$ and assume
$\alpha:S\rightarrow \mathbb{R}$ and $v:S\rightarrow \mathcal H$ satisfy
the so-called Glaeser--Whitney conditions:
\begin{equation}\label{eq:GW}
\begin{aligned}
&K_1:=\sup_{\substack{s_1,s_2\in S\\s_1\neq s_2}}
\frac{\left|
\alpha(s_2)-\alpha(s_1)
-\langle v(s_1),s_2-s_1\rangle
\right|}
{\|s_1-s_2\|^2}
<+\infty,\\[2ex]
&K_2:=\sup_{\substack{s_1,s_2\in S\\s_1\neq s_2}}
\frac{\|v(s_1)-v(s_2)\|}
{\|s_1-s_2\|}
<+\infty .
\end{aligned}
\end{equation}
\begin{theorem}\label{lem:whitney}(
  \(C^{1,1}\)-Glaeser--Whitney almost-minimal extension;
  see \cite[Theorem~3.1]{DaniilidisEtAl2018}
)
Let $S$ be a nonempty subset of a Hilbert space $H$ and let
$(\alpha(s),v(s))_{s\in S}$ be a $1$-Taylor field on $S$ satisfying
(\ref{eq:GW}). Then, the function
\[
G(x)=F(x)-\frac{1}{2}\overline{\mu}\|x\|^2
\]
is an explicit $C^{1,1}$-extension of the $1$-Taylor field $(\alpha,v)$,
provided that $F$ is the convex extension of the $1$-Taylor field
$(\tilde{\alpha},\tilde v)$ where for all $s\in S$,
$
\tilde{\alpha}(s):=\alpha(s)+\frac{1}{2}\overline{\mu}\|s\|^2
$
and
$
\tilde v(s):=v(s)+\overline{\mu}s,
$
with
$
\overline{\mu}
=
2K_1+K_2+
\sqrt{(2K_1+K_2)^2+K_2^2},
$
where $K_1,K_2$ are given by (\ref{eq:GW}).
For \(s_1\ne s_2\), define
\begin{align}
A_{s_1s_2}
&:=
\frac{
2\bigl(\alpha(s_1)-\alpha(s_2)\bigr)
+\ip{v(s_1)+v(s_2)}{s_2-s_1}}
{\norm{s_1-s_2}^{2}},\label{Afor}\\
B_{s_1s_2}
&:=
\frac{\norm{v(s_1)-v(s_2)}}{\norm{s_1-s_2}},\label{Bfor}
\end{align}
and set
\[
\Gamma_1(S,(\alpha,v))
:=
\sup_{s_1\ne s_2}
\left(
\sqrt{A_{s_1s_2}^{\,2}+B_{s_1s_2}^{\,2}}
+|A_{s_1s_2}|
\right).
\]
%For a globally defined \(C^{1,1}\) function \(h\),
%\[
% \Gamma_1(\Hsp,(h,\nabla h))=\Lip(\nabla h)
%\]
%\cite{LeGruyer}.
Moreover, the extension $G$, which satisfies
$G(s)=\alpha(s)$ and $\nabla G(s)=v(s)$ for every $s\in S$,
 is almost minimal, i.e.
\[
\Gamma_1(S,(\alpha,v))
\leq
\Gamma_1(H,(G,\nabla G))
=
\operatorname{Lip}(\nabla G)
\leq
\frac{5+\sqrt{29}}{2}
\Gamma_1(S,(\alpha,v)).
\]
\end{theorem}

%Let $E$ be a nonempty subset of a Hilbert space and let $a:E\to\R$, $v:E\to H$. For $x\neq y$, set
%\[
%A_{x,y}:=\frac{2(a(x)-a(y))+\ip{v(x)+v(y)}{y-x}}{\norm{x-y}^2},\qquad
%B_{x,y}:=\frac{\norm{v(x)-v(y)}}{\norm{x-y}},
%\]
%and
%\[
%\Gamma^1(E,(a,v)):=\sup_{x\neq y}\left(\sqrt{A_{x,y}^2+B_{x,y}^2}+|A_{x,y}|\right).
%\]
%
%We review the almost-minimal extension result; see, e.g., \cite[Theorem 3.1]{DaniilidisEtAl2018}.
%\begin{theorem}\label{thm:whitney}
%If the Glaeser--Whitney compatibility conditions hold on $E$, then there exists $G\in C^{1,1}(H)$ with $G=a$ and $\nabla G=v$ on $E$, and
%\[
%\Gamma^1(E,(a,v))\le\Lip(\nabla G)\le \frac{5+\sqrt{29}}{2}\,\Gamma^1(E,(a,v)).
%\]
%\end{theorem}


For the counterexample construction, we do not need an explicit formula for
the extension in Theorem~\ref{lem:whitney}. Instead, we shall invoke the
following  \(C^{1,1}\) Whitney extension criterion.



\begin{corollary}[A practical  \(C^{1,1}\) Whitney extension criterion]
\label{cor:practical-whitney}
Let \(E\subseteq\mathbb R^m\) be nonempty and let
\[
a:E\to\mathbb R,
\qquad
v:E\to\mathbb R^m.
\]
Suppose that there exists \(M>0\) such that, for every \(p,q\in E\),
\begin{align}
\|v(p)-v(q)\|
&\le M\|p-q\|,
\label{eq:whitney-gradient-condition}\\
\bigl|a(p)-a(q)-\langle v(q),p-q\rangle\bigr|
&\le M\|p-q\|^2.
\label{eq:whitney-value-condition}
\end{align}
Then there exists a function \(H\in C^{1,1}(\mathbb R^m)\) such that
\[
H|_E=a,
\qquad
\nabla H|_E=v,
\]
and
\[
\operatorname{Lip}(\nabla H)
\le C_{\mathrm W}M,
\]
where
\[
C_{\mathrm W}
:=
\frac{5+\sqrt{29}}{2}\bigl(3+\sqrt{10}\bigr)
<32.
\]
In particular, \(C_{\mathrm W}\) is a universal constant independent of
the dimension \(m\).
\end{corollary}

\begin{proof}
For distinct \(p,q\in E\), let \(A_{p,q}\) and \(B_{p,q}\) be defined as in
\eqref{Afor} and \eqref{Bfor}, respectively. By
(\ref{eq:whitney-gradient-condition}),
$
B_{p,q}\le M.
$
Define
\[
R_{p,q}
:=
a(p)-a(q)-\langle v(q),p-q\rangle.
\]
By \eqref{eq:whitney-value-condition},
$
|R_{p,q}|\le M\|p-q\|^2.
$
Moreover,
\begin{align*}
&2\bigl(a(p)-a(q)\bigr)
 +\langle v(p)+v(q),q-p\rangle\\
&\qquad
=2R_{p,q}-\langle v(p)-v(q),p-q\rangle.
\end{align*}
Consequently,
\[
|A_{p,q}|
\le
\frac{2|R_{p,q}|}{\|p-q\|^2}
+
\frac{\|v(p)-v(q)\|}{\|p-q\|}
\le 3M.
\]
Therefore,
\[
\Gamma_1(E,(a,v))
\le
\sqrt{(3M)^2+M^2}+3M
=
(3+\sqrt{10})M.
\]
The conclusion follows directly from Theorem~\ref{lem:whitney}.
\end{proof}


\section{Counterexamples}\label{sec:construction}
\subsection{Basic Construction}
\begin{theorem}\label{thm:main-counterexample}
There exists an instance of problem~\eqref{eq:fractional-problem}
satisfying Assumption~\ref{ass:standing} such that  ${\rank}(A)=1$ and $z^k\neq0$ for every $k\ge 0$; moreover, the S-FSPS iterates satisfy
\[
\clust(x^k)=\{1\}\times\Sph^1,
\qquad
\sum_{k=0}^{\infty}\norm{x^{k+1}-x^k}=+\infty.
\]
Moreover, every point in the cluster set is a limiting lifted stationary point.
\end{theorem}
\begin{proof}
\textbf{Step 1: problem data.}
Let $n=3$, $s=2$, $p=3$, and
\[
A=\begin{pmatrix}1&0&0\\0&0&0\end{pmatrix},\qquad
{\cal S}:=\{1\}\times \overline{\mathbb B}_{\R^2},\qquad K=I_3.
\]
Set
\[
g(y)=|y_1|+|y_2|,\qquad f(x)\equiv1.
\]
Then \(g^*=\iota_{[-1,1]^2}\) and
$
\prox_{g^*/\gamma}=\Proj_{[-1,1]^2}.$
Choose \(\chi=2\) and \(\gamma_k=1/(k+K_0)\). We shall construct \(h\) so that
\(\operatorname{Lip}(\nabla h)\le 1\). Since \(\sigma_A=1\), we may take \(1\)
as an upper bound for the Lipschitz constant of \(\nabla h\),
 not necessarily its
exact value. Consequently,
\[
\delta_k=2(k+K_0+1),
\qquad
\alpha_k:=\frac{1}{\delta_k}.
\]

%The construction has three components. First, we prescribe a spiral approaching the unit circle while accumulating an infinite angular displacement. Second, we define gradient and function-value data along the orbit so that a projected-gradient step reproduces the next spiral point. Third, we verify the Whitney compatibility conditions and extend these data to a globally $C^{1,1}$ function.

Before stating the main construction, we briefly describe the
auxiliary sequences. The sequence \(\{d_k\}_{k\ge 0}\) measures the radial
distance of the prescribed iterates from the limiting unit circle and is chosen
to vanish as \(k\to\infty\). The sequence \(\{\beta_k\}_{k\ge 0}\) specifies the angular
increment between successive iterates; its nonsummability ensures that the
trajectory continues to rotate indefinitely. These two sequences determine the spiral
\(\{p^k\}_{k\ge 0}\). Given this orbit, \(\{v_k\}_{k\ge 0}\) is chosen as the gradient value required
for  reproducing \(p^{k+1}\) exactly from
\(p^k\). Finally, \(\{a_k\}_{k\ge 0}\) prescribes the corresponding function values so that
the pairs \(\{(a_k,v_k)\}_{k\ge 0}\) form a compatible first-order Whitney jet.
Throughout the construction, $C>0$ denotes a constant independent of $k$, $K_0$ and $\eta$, whose value may change from line to line.


\medskip\noindent
\textbf{Step 2: the spiral.}
Let $\eta\in(0,1/2]$ be specified later, and choose an integer
$K_0:=K_0(\eta)$ sufficiently large such that
\begin{equation}\label{eq:K0-conditions}
K_0>e^4,\qquad
d_0:=\frac{1}{\sqrt{\log K_0}}\leq\eta^2,
\qquad
\eta\alpha_0d_0^2\leq1.
\end{equation}
For $k\geq0$, define
\begin{eqnarray}\label{dkvark}
&&d_k:=\frac1{\sqrt{\log(k+K_0)}},\qquad r_k:=1-d_k,\nn\\
&&\varphi_0=0,\qquad \beta_k:=\varphi_{k+1}-\varphi_k=\eta\alpha_kd_k^2,\end{eqnarray}
and
\begin{eqnarray}\label{uk}
u(\varphi):=(\cos\varphi,\sin\varphi),\qquad p^k:=r_ku(\varphi_k),\qquad x^k:=(1,p^k).
\end{eqnarray}
We have
\begin{align}
 &d_k\le\eta^2\;\mbox{ for all}\; k\ge0,\label{adddk}\\
&\eta\alpha_kd_k^2\le 1\;\mbox{ for all}\; k\ge 0,\nn\\
 &\sum_{k=0}^\infty\beta_k=\frac\eta2\sum_{k=0}^\infty\frac1{(k+K_0+1)\log(k+K_0)}=\infty,
\;
\mbox{whereas}\; \beta_k\to0.
\end{align}
 Fix \(\vartheta\in[0,2\pi)\).  For every sufficiently large integer \(m\),
let \(k_m\) be the first index such that
\(\varphi_{k_m}\geq2\pi m+\vartheta\).  Then
\[
    0\leq\varphi_{k_m}-(2\pi m+\vartheta)
    \leq\beta_{k_m-1}\longrightarrow0,
\]
so \(p^{k_m}\to(\cos\vartheta,\sin\vartheta)\).  Since \(r_k\to1\), every
cluster point of \(\{p^k\}\) lies on \(\Sone\).  Hence
\[
    \operatorname{cluster}(p^k)=\Sone,
    \qquad
    \operatorname{cluster}(x^k)=\{1\}\times\Sone.
\]
 In addition, we prove the chord estimate stated in Lemma \ref{lem:chord-lower-bound}, i.e.,
%\[
%\norm{p^{k+1}-p^k}^2=(r_{k+1}-r_k)^2+4r_kr_{k+1}\sin^2(\beta_k/2),
%\]
%so, using $r_k,r_{k+1}\ge1/2$ and $\sin t\ge2t/\pi$ on $[0,\pi/2]$,
\begin{equation}\label{eq:chord}
\norm{p^{k+1}-p^k}\ge\frac{\beta_k}{\pi}.
\end{equation}
Therefore $\sum_{k=0}^{+\infty}\norm{p^{k+1}-p^k}=+\infty$.

\medskip\noindent
\textbf{Step 3: prescribed jet and elementary estimates.}
Define
\begin{eqnarray}\label{vkak}
v_k:=\frac{p^k-p^{k+1}}{\alpha_k},\qquad a_k:=\sum_{\ell=k}^{\infty}\alpha_\ell\norm{v_\ell}^2.
\end{eqnarray}
Then
\[
p^{k+1}=p^k-\alpha_kv_k,\qquad a_{k+1}-a_k=-\alpha_k\norm{v_k}^2.
\]
%By \eqref{eq:a-size} of Lemma~\ref{lem:elementary}, we have
%\(a_k<+\infty\).
%Choose \(K_0\) as specified in Lemma~\ref{lem:elementary}, and set
%\[
%C:=\max\{\widehat C_1,\widehat C_2,\widehat C_3,\widehat C_4\}.
%\]
By \eqref{eq:a-size} of Lemma~\ref{lem:elementary}, we have $a_k<+\infty$.
Set
\[
    C:=\max\{\widehat C_1,\widehat C_2,
             \widehat C_3,\widehat C_4\}.
\]
Then
\begin{eqnarray*}
&&0<d_k-d_{k+1}\leq C\alpha_kd_k^3,\label{eq:d-est}\\
&&\norm{v_k}\leq C\eta d_k^2,\label{eq:v-est}\\
&&\norm{v_{k+1}-v_k}\leq C\eta\norm{p^{k+1}-p^k},\label{eq:vdiff-est}\\
&&0\leq a_k\leq C\eta^2d_k^2.\label{eq:a-est}
\end{eqnarray*}
%The only delicate point is \eqref{eq:vdiff-est}. The decomposition used in the original proof yields
%\[
%\norm{v_{k+1}-v_k}\le C\alpha_kd_k^3+C\eta\alpha_kd_k^4+C\eta^2\alpha_kd_k^4+\text{higher-order terms}.
%\]
%By the strengthened choice $d_k\le d_0\le\eta^2$, every term is bounded by $C\eta^2\alpha_kd_k^2$. Combining this with \eqref{eq:chord} yields \eqref{eq:vdiff-est}.

Let
\[
E:=\Sph^1\cup\{p^k:k\ge0\},
\]
and prescribe
\[
a(q)=0,\quad v(q)=0\quad(q\in\Sph^1),\qquad a(p^k)=a_k,\quad v(p^k)=v_k.
\]

\medskip\noindent
\textbf{Step 4: spiral--circle estimates.}
Fix \(k\geq0\) and \(q\in\Sone\), and set \(R=\norm{p^k-q}\).
Because \(\norm{p^k}=1-d_k\),
\[
   d_k= \dist(p^k,\Sone)\le R.
\]
Using \(d_k\leq1\), \(0<\eta\leq1\), and
Lemma~\ref{lem:elementary}, let
\[
{C_*}:=\max\{\widehat C_1,\widehat C_2,\widehat C_3,\widehat C_4\}.
\]
Then
\begin{align}
    \norm{v(p^k)-v(q)}
    &=
    \norm{v_k}
    \leq C_*\eta d_k^2
    \leq C_*\eta R,
    \label{eq:mixed-gradient}\\
    |a(p^k)-a(q)-\ip{v(q)}{p^k-q}|
    &=
    a_k
    \leq C_*\eta^2d_k^2
    \leq C_*\eta R^2,
    \label{eq:mixed-forward}\\
    |a(q)-a(p^k)-\ip{v(p^k)}{q-p^k}|
    &\leq
    a_k+\norm{v_k}R \nonumber\\
    &\leq
    C_*\eta^2d_k^2+C_*\eta d_k^2R \nonumber\\
    &\leq
    C_*\eta R^2+C_*\eta R^2
    \leq 2 C_*\eta R^2.
    \label{eq:mixed-reverse}
\end{align}
 The last inequalities follow from
\(0<\eta\leq1\), \(d_k\leq R\), and \(d_k\leq1\). Thus, the gradient condition \eqref{eq:whitney-gradient-condition} and both
orientations of the value condition \eqref{eq:whitney-value-condition} are
verified for every spiral--circle pair by taking $C\geq 2C_*.$


\medskip\noindent
\textbf{Step 5: local angular separation spiral--spiral estimates.}
Let $i>j$ and suppose that
\[
0\le \varphi_i-\varphi_j\le \pi.
\]
Define
\begin{eqnarray}\label{Lij}
L_{j,i}:=\sum_{\ell=j}^{i-1}\norm{p^{\ell+1}-p^\ell}.
\end{eqnarray}
Then, Proposition~\ref{prop:local-half-turn} gives
\begin{eqnarray}\label{add1}
\norm{v_i-v_j}\le \widehat C_5\eta\norm{p^i-p^j}.
\end{eqnarray}
Moreover, Proposition~\ref{prop:local-quadratic} yields
\begin{eqnarray}\label{add2}
\left|a_i-a_j-\ip{v_j}{p^i-p^j}\right|
\le
\widehat C_6\eta\norm{p^i-p^j}^2,
\end{eqnarray}
Enlarging $C$ further, take
$C\ge \max(\widehat C_5,\widehat C_6)$.


\medskip\noindent
\textbf{Step 6: large angular separation spiral--spiral estimates.}
If $i>j$ and $\varphi_i-\varphi_j\ge\pi$, then
Lemma~\ref{lem:separation} yields
\[
\norm{p^i-p^j}
\ge
c\max\{d_i,d_j\},
\qquad
c:=1-e^{-\pi/2}>0.
\]
Next, we prove the following estimates.
\begin{eqnarray}\label{eq:different-gradient}
\norm{v_i-v_j}
\le
C\eta\norm{p^i-p^j},
\end{eqnarray}
and
\begin{eqnarray}\label{eq:different-value}
\left|a_i-a_j-\ip{v_j}{p^i-p^j}\right|
\le
C\eta\norm{p^i-p^j}^2.
\end{eqnarray}
First, let
\[
R:=\norm{p^i-p^j},
\qquad
D:=\max\{d_i,d_j\}=d_j.
\]
By Lemma~\ref{lem:separation},
$D\le \frac{R}{c}.$
Since $D\le1$, Lemma~\ref{lem:elementary} gives
\begin{align*}
&\norm{v_i-v_j}
\le
\norm{v_i}+\norm{v_j}
\le
2{\widehat C_2}\eta D^2
\le
2{\widehat C_2}\eta R/c,
\\
&\left|a_i-a_j-\ip{v_j}{p^i-p^j}\right|
\le
a_i+a_j+\norm{v_j}R
\notag\\
&\le
2\widehat C_4\eta^2D^2+\widehat C_2\eta D^2R
\notag\le
\left(\frac{2\widehat C_4}{c^2}+\frac{\widehat C_2}{c}\right)\,\eta R^2,
\end{align*}
where we used \(0<\eta\leq1\). Enlarge $C$ further so that
\[
    C\geq
    \max\left\{
        \frac{2\widehat C_2}{c},
        \frac{2\widehat C_4}{c^2}
        +\frac{\widehat C_2}{c}
    \right\}.
\]
Then \eqref{eq:different-gradient} and \eqref{eq:different-value} hold.


\medskip\noindent
\textbf{Step 7: Whitney extension and exact realization.}
The circle--circle estimates are trivial because \(a=0\) and \(v=0\) on
\(\Sone\).  Equations
\eqref{eq:mixed-gradient}--\eqref{eq:mixed-reverse} treat all mixed pairs.
Equations \eqref{add1}--\eqref{add2} treat
spiral--spiral pairs with angular separation at most \(\pi\), while
\eqref{eq:different-gradient}--\eqref{eq:different-value} treat the remaining
pairs.
The displayed value estimates were written with the earlier spiral point as
the base point.  The reverse orientation follows from
\begin{eqnarray}
    &&a_j-a_i-\ip{v_i}{p^j-p^i}\nonumber\\
    &&=
    -\left(a_i-a_j-\ip{v_j}{p^i-p^j}\right)
    +\ip{v_i-v_j}{p^i-p^j}.
\end{eqnarray}
Together with the corresponding gradient estimates, see
\eqref{eq:mixed-gradient}, \eqref{eq:different-gradient}, and
\eqref{add1},  the preceding estimates show that, after enlarging
\(C>0\) if necessary, the jet \((a,v)\) satisfies
\[
    \|v(p)-v(q)\|\le C\eta\|p-q\|,
\]
and
\[
    |a(p)-a(q)-\langle v(q),p-q\rangle|
    \le C\eta\|p-q\|^2
\]
for all \(p,q\in E\).
Therefore, Corollary~\ref{cor:practical-whitney} yields that there exists
\(\widetilde h_0\in C^{1,1}(\mathbb R^2)\) such that
\[
\widetilde h_0(p^k)=a_k,
\qquad
\nabla\widetilde h_0(p^k)=v_k,
\]
and
\[
\widetilde h_0|_{\mathbb S^1}=0,
\qquad
\nabla\widetilde h_0|_{\mathbb S^1}=0,
\qquad
\operatorname{Lip}(\nabla\widetilde h_0)
\le C_{\mathrm W}C\eta.
\]
The preceding estimates yield a constant \(C>0\), independent of
\(\eta\), \(K_0\), and \(k\), such that the two Whitney bounds above
hold for all \(p,q\in E\). Choose this final constant \(C\) large
enough to dominate all the constants appearing in Steps~3--7.
We then choose
\[
    0<\eta\leq
    \min\left\{\frac12,\frac{1}{C_WC}\right\}.
\] 
Here, $C\ge2\max\{2 C_*,\widehat C_5,\widehat C_6,\frac{2\widehat C_2}{c},
\frac{2\widehat C_4}{c^2}+\frac{\widehat C_2}{c} \}$.
The initial choice of \(K_0=K_0(\eta)\) can now be made sufficiently large so that all the preceding elementary estimates and condition (\ref{eq:K0-conditions}) hold.
%Then choose $K_0=K_0(\eta)$ sufficiently large so that all elementary estimates and condition (\ref{eq:K0-conditions}) hold.
Consequently,
\[
\operatorname{Lip}(\nabla\widetilde h_0)\le1.
\]


%Together with the corresponding gradient estimates; see
%\eqref{eq:mixed-gradient}, \eqref{eq:different-gradient}, and
%\eqref{add1}, we conclude that
%Circle--circle pairs are trivial. Steps 4--6 yield the Glaeser--Whitney compatibility estimates on all of $E$, hence
%\[
%\Gamma^1(E,(a,v))\le C\eta.
%\]
%Choose $\eta>0$ so small that
%\begin{equation}\label{eq:eta}
%\frac{5+\sqrt{29}}{2}C\eta\le1.
%\end{equation}
%Theorem~\ref{thm:whitney} yields $\widetilde h_0\in C^{1,1}(\R^2)$ satisfying
%\[
%\widetilde h_0(p^k)=a_k,\qquad \nabla\widetilde h_0(p^k)=v_k,
%\]
%\[
%\widetilde h_0|_{\Sph^1}=0,\qquad \nabla\widetilde h_0|_{\Sph^1}=0,
%\qquad \Lip(\nabla\widetilde h_0)\le1.
%\]
%Circle--circle pairs are trivial, while Steps~4--6 show that there exists
%an absolute constant \(C>0\) such that, for every \(p,q\in E\),
%\[
%\|v(p)-v(q)\|\le C\eta\|p-q\|
%\]
%and
%\[
%\bigl|a(p)-a(q)-\langle v(q),p-q\rangle\bigr|
%\le C\eta\|p-q\|^2.
%\]
%By Corollary~\ref{cor:practical-whitney},
Add a constant $c_0$ so that $\widetilde h:=\widetilde h_0+c_0\ge1$ on $\overline{\mathbb B}_{\R^2}$, and define
\[
h(x_1,x_2,x_3):=\widetilde h(x_2,x_3).
\]
Initialize $x^0=(1,p^0)$, $z^0=(1,0)$, and $\theta_0>0$. Since $Ax=(1,0)$ for all $x\in\cal S$,
\[
z^{k+1}=\Proj_{[-1,1]^2}\left(\frac1{\gamma_k},0\right)=(1,0).
\]
Together with the initialization $z^0=(1,0)$, this yields  $z^k=(1,0)\neq0$ for all $k\ge0$. Also
\[
A^*z^k=(1,0,0),\qquad \nabla h(x^k)=(0,v_k),\ f\equiv 1,\; y^{k+1}=0.
\]
Therefore, the term  $\frac{\theta_k}{\delta_k}K^* y^{k+1}$ vanishes.
\[
x^k-\frac1{\delta_k}\nabla h(x^k)-\frac1{\delta_k}A^*z^k=(1-\alpha_k,p^{k+1}).
\]
Since $\Proj_{\cal S}(t,p)=(1,\Proj_{\overline{\mathbb B}_{\R^2}}p)$ and $p^{k+1}\in \overline{\mathbb B}_{\R^2}$, the $x$-update gives
\[
x^{k+1}=(1,p^{k+1}).
\]
Thus the prescribed spiral is exactly the primal S-FSPS orbit.

It remains to  verify Assumption \ref{ass:standing}.
\[
A({\cal S})=\{(1,0)\}\subseteq\dom\partial g,\qquad \partial g(1,0)=\{ 1 \}\times[-1,1],
\]
so $\dist(0,\partial g(Ax))=1$ for every $x\in{\cal S}$. Moreover,
\[
g(Ax)+h(x)=1+h(x)\ge2.
\]
Hence Assumption \ref{ass:standing} holds and the cluster-set and infinite-length conclusions now follow from Step 2.

We conclude by verifying that every point in the cluster set is a limiting lifted stationary point.
Let $\overline x=(1,q)$ where $q\in\Sone$. Then,
\begin{eqnarray}\nabla h(\overline x)=0,\;\partial g(A\overline x)=\{1\}\times [-1,1].\end{eqnarray}
$\overline z=(1,0)\in \partial g(A\overline x)$.
Since $(-1,0,0)\in N_{\cal S}(\overline x)$, we have
\begin{eqnarray*}0\in A^*\partial g(A{\overline x})+\nabla h(\overline x)+N_{\cal S}(\overline x).\end{eqnarray*}
Because $f\equiv 1$, this implies that $\overline x$ is  a limiting lifted stationary point.
\end{proof}

%\begin{remark}
%In the construction used in the proof of
%Theorem~\ref{thm:main-counterexample}, the set \(A(\mathcal S)\) is a
%singleton, and hence \(g\circ A\) is constant on \(\mathcal S\). This feature
%does not diminish the logical force of the counterexample, since a single
%admissible instance suffices to refute a general claim of whole-sequence
%convergence. %Nevertheless, the example is not entirely trivial:
%%\(A\neq 0\) and \(z^k\neq 0\) for every \(k\), so all the dual iterates are
%%nonzero.
%\end{remark}

\begin{remark}
The counterexample does not contradict the subsequential convergence theorem in
\cite{BotLiTao2025}. It shows that different subsequences may converge to
different stationary cluster points. Moreover, the constructed instance satisfies
\[
    \gamma_k\downarrow0,\qquad
    \sum_{k=0}^{\infty}\gamma_k=+\infty,
    \qquad
    \sum_{k=0}^{\infty}\gamma_k^2<+\infty.
\]
Thus, even square summability of the smoothing parameters does not
ensure whole-sequence convergence.
\end{remark}
\subsection{Full-Dimensional Fractional Construction}
We next construct a full-dimensional counterexample for the fractional model satisfying several natural nondegeneracy properties.

%\begin{theorem}\label{thm:main-counterexample:SCP}
%Under Assumption~\ref{ass:standing}, whole-sequence convergence of S-FSPS cannot in general be guaranteed. More precisely, there exists a structured fractional program satisfying all the standing assumptions for which
% \(\mathcal S\) is full-dimensional,  with full-row-rank $A$, nonconstant $f\circ K$, and nonconstant $g\circ A$
% and $z^k\neq0$ for every $k$, while
%\[
%\clust(x^k)=\{1\}\times\Sph^1,
%\qquad
%\sum_{k=0}^{\infty}\norm{x^{k+1}-x^k}=+\infty.
%\]
%\end{theorem}
\begin{theorem}\label{thm:full-dimensional-counterexample}
There exists an instance of problem~\eqref{eq:fractional-problem}
satisfying Assumption~\ref{ass:standing} such that
\(\operatorname{int}\mathcal S\neq\varnothing\), \(A\) has full row
rank, and each of \(f\circ K\), \(g\circ A\), and \(g\circ A+h\) is
nonconstant on \(\mathcal S\). Moreover, the S-FSPS iterates satisfy
\[
    \operatorname{cluster}(x^k)=\{1\}\times\mathbb S^1,
    \qquad
    \sum_{k=0}^{\infty}\|x^{k+1}-x^k\|=+\infty.
\]
Furthermore, every point in the cluster set is a limiting lifted stationary point.
Consequently, since the cluster set contains infinitely many points,
the primal sequence does not converge.
\end{theorem}

%The proof should verify the following properties separately:
%\[
%    K(\mathcal S)=[1,2]
%    \subseteq \operatorname{int}(\operatorname{dom}f),
%    \qquad
%    f(Kx)=3-x_1\in[1,2],
%\]
%\[
%    \operatorname{rank}A=2,
%    \qquad
%    \|A\|=1,
%\]
%\[
%    A(\mathcal S)\subset(0,+\infty)^2,
%    \qquad
%    \partial g(Ax)=\{(1,1)\},
%\]
%\[
%    \operatorname{dist}\bigl(0,\partial g(Ax)\bigr)=\sqrt{2},
%\]
%and
%\[
%    g(Ax)+h(x)=\widehat h(x_2,x_3)\geq 3.
%\]
%
%Note that the subdifferential is a set; hence one should write
%\[
%    \partial g(w)=\{(1,1)\},
%\]
%rather than $\partial g(w)=(1,1)$.

\begin{proof}
Let
\[
\mathcal S:=[1,2]\times \overline{\mathbb B}_{\R^2},
\qquad
g(u):=\|u\|_1 \quad (u\in\mathbb R^2),
\qquad
f(s):=3-s \quad (s\in\mathbb R),
\]
and define
\[
A=\frac13
\begin{pmatrix}
2&1&0\\
2&0&1
\end{pmatrix},
\qquad
K=
\begin{pmatrix}
1&0&0
\end{pmatrix},
\qquad
b=
\begin{pmatrix}
4/3\\
1/3\\
1/3
\end{pmatrix}.
\]
Then
$\sigma_A=1$ and $\operatorname{rank}(A)=2.$
Moreover, for every \(x\in\mathcal S\),
$
Ax\geq
\begin{pmatrix}
1/3\\
1/3
\end{pmatrix},$
$$g(Ax)=\|Ax\|_1=\frac{4}{3}x_1+\frac{1}{3}x_2+\frac{1}{3}x_3=\langle b,x\rangle.$$

\noindent Choose \(\chi=2\) and set
$
\gamma_k:=\frac{1}{k+K_0}.
$
We shall construct a function \(h\) satisfying
\[
\operatorname{Lip}(\nabla h)\leq 1.
\]
Accordingly, we set
$
\delta_k:=2(k+K_0+1)$ where $K_0\ge e^4>4$
and
$\alpha_k:=\frac{1}{\delta_k}.$
%Choose $\eta\in(0,1/2]$ and then choose $K_0=K_0(\eta)$
%exactly as in Step~7 of the proof of Theorem \ref{thm:main-counterexample}; in particular,
%\eqref{eq:K0-conditions} holds.
Choose $\eta>0$ and $K_0=K_0(\eta)$ exactly as in Step~7 of the
proof of Theorem~\ref{thm:main-counterexample}; in particular,
\[
    0<\eta\leq
    \min\left\{\frac12,\frac{1}{C_WC}\right\},
\]
and $K_0$ is sufficiently large so that \eqref{eq:K0-conditions}
holds.
Then, define \(d_k\), \(r_k\), \(\beta_k\), \(p^k\),
and \(x^k\) as in \eqref{dkvark} and \eqref{uk}, respectively.
Moreover, let \(a_k\) and \(v_k\) be defined as in \eqref{vkak}.

Following Steps 2--6 and the Whitney-extension part of Step 7 in the
proof of Theorem 4.1, and applying Corollary 3.2, we obtain
\(h_0\in C^{1,1}(\mathbb R^2)\) such that
\[
h_0(p^k)=a_k,
\qquad
\nabla h_0(p^k)=v_k,
\]
and
$h_0|_{\Sph^1}=0$,
$
\nabla h_0|_{\Sph^1}=0$,
$\operatorname{Lip}(\nabla h_0)\leq 1.
$
%Choose \(c_0>0\) sufficiently large so that
%\[
%    \widehat h(u):=h_0(u)+c_0\geq 3
%    \qquad\text{for all }u\in\mathbb B_{\mathbb R^2}.
%\]
%and define
%$
%h(x_1,x_2,x_3)
%:=
%\widehat h(x_2,x_3)-\langle b,x\rangle.
%$
Choose $c_0>0$ sufficiently large so that
\[
    \widehat h(u):=h_0(u)+c_0\geq 3
    \qquad\text{for all }u\in\mathbb B_{\mathbb R^2}.
\]
Define
\[
    h(x_1,x_2,x_3)
    :=\widehat h(x_2,x_3)-\langle b,x\rangle.
\]
Thus, $\operatorname{Lip}(\nabla h)\leq 1$.
Initialize $x^0=(1,p^0)$, $z^0=(1,1)$, and $\theta_0>0$. Since $\gamma_k<1/3$ and $Ax^{k+1}\ge(1/3,1/3)$, it follows that
$z^{k+1}=\Proj_{[-1,1]^2}\left(\frac{A x^{k+1}}{\gamma_k}\right)=(1,1).$
Thus, $z^k=(1,1)\neq0$ for all $k\ge 0$. Also
\[
A^*z^k=b,\qquad \nabla h(x^k)=(0,v_k)-b,\ y^{k+1}=-1.
\]
Therefore,
\[
x^k+\frac{\theta_k}{\delta_k}K^* y^{k+1}-\frac1{\delta_k}\nabla h(x^k)-\frac1{\delta_k}A^*z^k=(1-\alpha_k\theta_k,p^{k+1}).
\]

We prove by induction that, for every $k\ge0$, the following two statements hold:
\begin{itemize}
\item[$\large{\textcircled{\small{1}}}_{k}$] $x^k=(1,p^k)$;
\item[$\large{\textcircled{\small{2}}}_{k}$] $\theta_k>0$.
\end{itemize}
The base case follows from the initialization $x^0=(1,p^0)$ and $\theta_0>0$.
Suppose that $\large{\textcircled{\small{1}}}_{k}$ and $\large{\textcircled{\small{2}}}_{k}$ hold.
Note that $x^{k+1}=(\Proj_{[1,2]}(1-\alpha_k\theta_k),\Proj_{\overline{\mathbb B}_{\R^2}}(p^{k+1}))=(1,p^{k+1})$
because $\theta_k>0$ and $p^{k+1}\in\overline{\mathbb B}_{\R^2}$.
Thus, $\large{\textcircled{\small{1}}}_{k+1}$ holds.

Moreover, since
\(g(Ax)=\langle b,x\rangle\) and
\(h(x)=\widehat h(x_2,x_3)-\langle b,x\rangle\), we have
$
\Psi(x^{k+1},z^{k+1};\gamma_k)
=\widehat h(p^{k+1})-\gamma_k>0,
\;
f(Kx^{k+1})=2.$ Hence,
$
\theta_{k+1}
=\frac{\widehat h(p^{k+1})-\gamma_k}{2}>0,
$
so \(\textcircled{\small 2}_{k+1}\) holds.
Thus the prescribed spiral is exactly the primal S-FSPS orbit.

It remains to  verify Assumption \ref{ass:standing}.
   $ K(\mathcal S)=[1,2]
    \subseteq \operatorname{int}(\operatorname{dom}f),$
    $f(Kx)=3-x_1\in[1,2],$
$
    \operatorname{rank}(A)=2,
    \;
    \|A\|=1,$
$A(\mathcal S)\subseteq
\{u\in\mathbb R^2:u_1>0,\ u_2>0\}.$
   For any $x\in{\cal S}$,
    $\partial g(Ax)=\{(1,1)^\top\}$.
Hence, $\dist(0,\partial g(Ax))=\sqrt{2}$ for every $x\in{\cal S}$. Moreover,
\[
g(Ax)+h(x)=\widehat h(x_2,x_3)\ge 3.
\]
Hence, all the standing assumptions are satisfied.

We conclude by verifying that every point in the cluster set is a limiting lifted stationary point.
Let \(\overline x=(1,q)\), where \(q\in\mathbb S^1\), and set
\(e_1:=(1,0,0)^\top\).

Since
$
\partial g(A\overline x)=\{(1,1)^\top\},\;
A^*(1,1)^\top=b,\;
\nabla h(\overline x)=-b,
$
we have
$
A^*(1,1)^\top+\nabla h(\overline x)=0.
$
Moreover,
$
f(K\overline x)=2,\;
g(A\overline x)+h(\overline x)=c_0,\;
K^*\partial f(K\overline x)=\{-e_1\}.
$
Because \(\overline x_1=1\), we have
$
-\frac{c_0}{2}e_1\in N_{\mathcal S}(\overline x).
$
Consequently,
$$
2\left(
A^*(1,1)^\top+\nabla h(\overline x)-\frac{c_0}{2}e_1
\right)
-c_0(-e_1)=0.
$$
Therefore,
$
0\in
f(K\overline x)\bigl(
A^*\partial g(A\overline x)+\nabla h(\overline x)+N_{\mathcal S}(\overline x)
\bigr)
-\bigl(g(A\overline x)+h(\overline x)\bigr)K^*\partial f(K\overline x),
$
so \(\overline x\) is a limiting lifted stationary point.
Moreover, the numerator is nonconstant on $\mathcal S$. Indeed,
$g(Ax)+h(x)=\widehat h(x_2,x_3).$
For every $q\in\mathbb S^1$, we have $\widehat h(q)=c_0$, whereas
$
\widehat h(p^k)=c_0+a_k>c_0,
$
where we have used $a_k>0$ for $k\geq 0$. Indeed,
$
r_{k+1}\neq r_k
\;\Longrightarrow\;
p^{k+1}\neq p^k
\;\Longrightarrow\;
v_k\neq 0,
$
and hence $a_k>0$ by \eqref{vkak}. Therefore, $g\circ A +h $ is nonconstant on \(\mathcal S\).
Since \(f(Kx)=3-x_1\) and \(g(Ax)=\langle b,x\rangle\),
both \(f\circ K\) and \(g\circ A\) are nonconstant on \(\mathcal S\).
\end{proof}
\section{Nonconvex  Composite Optimization Specialization}\label{sec:com}

%Our construction shows that the subsequential convergence result in
%\cite[Theorem~4.3]{BotLiTao2025} is essentially sharp under the standing
%assumptions. Even when the primal feasible set is compact, the
%smoothing sequence vanishes but is nonsummable, the linear operator appearing
%in the composite term is nonzero, and the dual iterates remain nonzero, the
%primal sequence may have a continuum of cluster points and trace an
%infinite-length orbit. Consequently, the convergence of the entire sequence cannot be ensured.

The same phenomenon  occurs for the  composite
problem.
\begin{equation}\label{eq:PP}
    \min_{x\in\mathcal S}\; g(Ax)+h(x),
\end{equation}
where $\mathcal S\subseteq\mathbb R^n$ is nonempty, convex and compact;
$g:\mathbb R^s\to\overline{\mathbb R}$ is proper, convex and lower
semicontinuous; $A:\mathbb R^n\to\mathbb R^s$ is linear; and
$h:\mathbb R^n\to\mathbb R$ is differentiable on an open set containing
$\mathcal S$, with a Lipschitz continuous gradient. Choose $\chi>1$, let $\{\gamma_k\}_{k\ge0}$ be positive and
nonincreasing and satisfy~\eqref{eq:gammak}, and define
$\delta_k$ by~\eqref{eq:delta-general}.
Starting from $x^0\in\mathcal S$ and $z^0\in\mathbb R^s$,
generate, for $k\geq 0$,
\begin{align}
x^{k+1}
&=
\Proj_{\mathcal S}\left(
x^k-\frac{1}{\delta_k}
\bigl(\nabla h(x^k)+A^*z^k\bigr)
\right),
\label{eq:reduced-x}\\
z^{k+1}
&=
\prox_{g^*/\gamma_k}
\left(\frac{Ax^{k+1}}{\gamma_k}\right).
\label{eq:reduced-z}
\end{align}
%where $\gamma_k$ and $\delta_k$ are defined in
%\eqref{eq:gammak} and \eqref{eq:delta-general}, respectively.

The following result establishes a subsequential convergence guarantee for
this reduced scheme.

\begin{theorem}\label{cor:reduced-subsequential}
Assume that $\mathcal S$ is nonempty,  convex, and compact; that
$h$ is differentiable on an open neighborhood of $\mathcal S$ and
$\nabla h$ is $L_{\nabla h}$-Lipschitz continuous there for some
$L_{\nabla h}>0$; that $g:\mathbb R^s\to\overline{\mathbb R}$ is proper, convex, and lower
semicontinuous; that
$A(\mathcal S)\subseteq\dom(\partial g)$; and that there exists
$\ell>0$ such that
$
    \dist\bigl(0,\partial g(Ax)\bigr)\leq\ell
    \;\text{for every }x\in\mathcal S.
$
Then every sequence $\{(x^k,z^k)\}_{k\ge 0}$ generated by
\eqref{eq:reduced-x}--\eqref{eq:reduced-z} satisfies
\begin{equation}\label{eq:reduced-liminf}
    \liminf_{k\to\infty}
    \delta_k\|x^{k+1}-x^k\|=0.
\end{equation}
Moreover, there exist a strictly increasing sequence of indices
\(\{k_j\}_{j\geq1}\) and a pair
\((\bar x,\bar z)\in\mathcal S\times\mathbb R^s\) such that
$
    x^{k_j}\to\bar x,\; z^{k_j}\to\bar z,
$
 and
\begin{equation}\label{eq:reduced-lifted-stationarity}
    \overline z\in\partial g(A\overline x),
    \qquad
    0\in
    \nabla h(\overline x)+A^*\overline z+N_{\mathcal S}(\overline x).
\end{equation}
Consequently, $\overline x$ is an exact stationary point of \eqref{eq:PP}, in
the sense that
\begin{equation}\label{eq:reduced-stationarity}
    0\in
    \nabla h(\overline x)
    +A^*\partial g(A\overline x)
    +N_{\mathcal S}(\overline x).
\end{equation}

\end{theorem}

\begin{proof}
%For $u\in A(\mathcal S)$ and $\gamma>0$, Moreau's decomposition gives
%\[
%    \prox_{g^*/\gamma}\left(\frac{u}{\gamma}\right)
%    =
%    \frac{1}{\gamma}
%    \bigl(u-\prox_{\gamma g}(u)\bigr).
%\]
%Let
%\[
%    p=\prox_{\gamma g}(u)
%    \quad\text{and}\quad
%    w=\frac{u-p}{\gamma}.
%\]
%Then $w\in\partial g(p)$. For any $v\in\partial g(u)$, the monotonicity of
%$\partial g$ yields
%\[
%    \langle w-v,p-u\rangle\geq0.
%\]
%Since $p-u=-\gamma w$, it follows that
%\[
%    \|w\|^2\leq\langle v,w\rangle
%    \leq\|v\|\,\|w\|,
%\]
%and hence $\|w\|\leq\|v\|$. Taking the infimum over
%$v\in\partial g(u)$ gives
%\[
%    \left\|
%    \prox_{g^*/\gamma}\left(\frac{u}{\gamma}\right)
%    \right\|
%    \leq
%    \dist\bigl(0,\partial g(u)\bigr).
%\]
%Consequently, the dual update satisfies
Invoking Lemma~\ref{lem:uniform-moreau-estimates}(i) together with the
dual update~\eqref{eq:reduced-z}, we obtain
\begin{eqnarray}\label{zkbound}
    \|z^{k+1}\|
    =\|\nabla g_{\gamma_k}(A x^{k+1})\|\le \ell.
\end{eqnarray}
Hence, the dual sequence \(\{z^k\}\) is bounded.
Moreover,
$
    A(\mathcal S)\subseteq\dom(\partial g)\subseteq\dom g.
$
Thus $g\circ A+h$ is finite-valued and lower semicontinuous on the compact
set $\mathcal S$, and therefore
$
    \inf_{x\in\mathcal S}\bigl(g(Ax)+h(x)\bigr)>-\infty.
$

Define $\Psi_k:=\Psi(x^k,z^k;\gamma_{k-1})$ for $k\ge 1$. By the dual update (\ref{eq:reduced-z}) at \((k-1)\)th iteration,
\begin{eqnarray}\label{zkop}
    z^k=\nabla g_{\gamma_{k-1}}(Ax^k),
\end{eqnarray}
and hence
\[
    \Psi_k=g_{\gamma_{k-1}}(Ax^k)+h(x^k).
\]
%As in the proof of \cite[Theorem~4.3(i)]{BotLiTao2025}, we have
%$$\Psi_{k+1}
%    \leq
%    \Psi_k
%    -{\widetilde c}_k\left\|x^{k+1}-x^k\right\|^2
%    + \Xi^{k+1},$$

Let $d^k:=x^{k+1}-x^k$ and define
$
    F_k:=g_{\gamma_{k-1}}\circ A+h,
    \;
    L_k:=L_{\nabla h}+\frac{\sigma_A^2}{\gamma_{k-1}}.
$
Since
the update \eqref{eq:reduced-x} is a
projected-gradient step for $F_k$ with stepsize $1/\delta_k$,
 the descent lemma yields
\[
    F_k(x^{k+1})
    \leq
    F_k(x^k)
    -
    \frac{\delta_k-L_k}{2}\|d^k\|^2.
\]
Since $\gamma_k\leq\gamma_{k-1}$, we have
\[
\begin{aligned}
    \delta_k-L_k
    =
    \chi\left(
        L_{\nabla h}+\frac{\sigma_A^2}{\gamma_k}
    \right)
    -
    \left(
        L_{\nabla h}+\frac{\sigma_A^2}{\gamma_{k-1}}
    \right)                                                  
    \geq
    (\chi-1)
    \left(
        L_{\nabla h}+\frac{\sigma_A^2}{\gamma_k}
    \right).
\end{aligned}
\]
Consequently,
\[
    g_{\gamma_{k-1}}(Ax^{k+1})+h(x^{k+1})
    \leq
    g_{\gamma_{k-1}}(Ax^k)+h(x^k)
    -
    \widetilde c_k\|d^k\|^2,
\]
where
\[
    \widetilde c_k
    :=
    \frac{\chi-1}{2}
    \left(
        L_{\nabla h}+\frac{\sigma_A^2}{\gamma_k}
    \right)
    =
    \frac{\chi-1}{2\chi}\delta_k
    >0.
\]
Moreover, the dual representation of the Moreau envelope and the
definition of $z^{k+1}$ imply
\[
    g_{\gamma_k}(Ax^{k+1})
    \leq
    g_{\gamma_{k-1}}(Ax^{k+1})
    +
    \frac{\gamma_{k-1}-\gamma_k}{2}
    \|z^{k+1}\|^2.
\]
Combining the preceding two inequalities gives
\begin{eqnarray}\label{despsi}
    \Psi_{k+1}
    \leq
    \Psi_k-\widetilde c_k\|x^{k+1}-x^k\|^2
    +\Xi^{k+1},
\end{eqnarray}
 where
\begin{eqnarray*}
\Xi^{k+1} := \frac{\gamma_{k-1}-\gamma_{k}}{2}\|{z}^{k+1}\|^2 \geq 0.
\end{eqnarray*}
\noindent Lemma~\ref{lem:uniform-moreau-estimates}(iii) gives the lower bound
\[
    g_{\gamma_{k-1}}\!\left(Ax^k\right)
    \geq
    g\!\left(Ax^k\right)
    -
    \frac{\gamma_{k-1}\ell^2}{2},
\]
so $\{\Psi_k\}_{k\ge 1}$ is bounded below. Combining (\ref{zkbound}), (\ref{despsi}), and the fact that
$\gamma_{k-1}\geq\gamma_k$ yields
$$
    \sum_{k=1}^{\infty} \Xi^{k+1}
    \leq
    \frac{\ell^2}{2}
    \sum_{k=1}^{\infty}\left(\gamma_{k-1}-\gamma_k\right)
    <+\infty.
$$

Consequently,
\begin{eqnarray}\label{addinf}\sum_{k=1}^{\infty}
    \delta_k\left\|x^{k+1}-x^k\right\|^2
    <+\infty.\end{eqnarray}
Since \(L_{\nabla h}>0\),
$\frac{1}{\delta_k} =
    \frac{\gamma_k}
    {\chi\left(L_{\nabla h}\gamma_k+\sigma_A^2\right)}
    \geq
    \frac{\gamma_k}
    {\chi\left(L_{\nabla h}\gamma_0+\sigma_A^2\right)}.$
    Thus,
$\sum_{k=1}^{\infty}\gamma_{k}=+\infty
    \;\Longrightarrow\;
    \sum_{k=1}^{\infty}\frac{1}{\delta_k}=+\infty.$
It then follows that
$$
    \liminf_{k\to+\infty}
    \delta_k\left\|x^{k+1}-x^k\right\|
    =0.
$$
Suppose not. Then there exist $\varepsilon>0$ and $\widehat K$ such that
$\delta_k\|x^{k+1}-x^k\|\ge\varepsilon$ for all $k\ge\widehat K$. Hence
$
  \delta_k^2\,\|x^{k+1}-x^k\|^2 \;\ge\; \varepsilon^2$
  for all $k\ge\widehat K,$
and summing over $k\ge\widehat K$ gives
\begin{eqnarray*}
  \sum_{k\ge\widehat K}\frac{[\delta_k^2\,\|x^{k+1}-x^k\|^2]}{\delta_k}
  \;\ge\; \varepsilon^2\sum_{k\ge\widehat K}\frac{1}{\delta_k} \;=\; +\infty,
\end{eqnarray*}
which contradicts (\ref{addinf}).

\noindent Choose a subsequence $\{k_j\}_{j\in\mathbb N}$ such that
$ \delta_{k_j}\|x^{k_j+1}-x^{k_j}\|\to0.$
By the compactness of $\mathcal S$ and the boundedness of the dual sequence,
we may pass to a further subsequence, without relabeling, such that
$x^{k_j}\to\overline x,\; z^{k_j}\to\overline z.$
Since $\delta_k$ is bounded away from zero under
\eqref{eq:delta-general}, we also have
$x^{k_j+1}-x^{k_j}\to0
    \;\text{and hence}\;
    x^{k_j+1}\to\overline x.$
The optimality condition for the primal update is
\begin{equation}\label{eq:reduced-x-optimality}
   \xi_j\in N_{\mathcal S}(x^{k_j+1}),
\end{equation}
where $\xi_j:=-[ \nabla h(x^{k_j})
    +A^*z^{k_j}
    +\delta_{k_j}(x^{k_j+1}-x^{k_j})]$.
On the other hand, the dual update at iteration $k_j-1$ gives
\begin{equation}\label{eq:reduced-z-optimality}
    Ax^{k_j}-\gamma_{k_j-1}z^{k_j}
    \in\partial g^*(z^{k_j}).
\end{equation}
Passing to the limit in \eqref{eq:reduced-x-optimality} and
\eqref{eq:reduced-z-optimality}, using $\gamma_{k_j-1}\to0$, $\xi_j\to -[\nabla h(\overline x)+A^*\overline z]$ and the closedness of
the graphs of $N_{\mathcal S}$ and $\partial g^*$ \cite[Proposition~8.7]{RockWets}, yields
$
    0\in
    \nabla h(\overline x)+A^*\overline z+N_{\mathcal S}(\overline x)
$
and
$
    A\overline x\in\partial g^*(\overline z).
$
The latter relation is equivalent to
$\overline z\in\partial g(A\overline x)$. This proves
\eqref{eq:reduced-lifted-stationarity}, and
\eqref{eq:reduced-stationarity} follows immediately.
\end{proof}
\begin{corollary}\label{cor:counterexample-consequences}

Apply \eqref{eq:reduced-x}--\eqref{eq:reduced-z}  to the instance constructed in Theorem
\ref{thm:main-counterexample}, with the same choices
\[
    \chi=2,\qquad
    \gamma_k=\frac{1}{k+K_0},\qquad
    \delta_k=2(k+K_0+1),
\]
and initialize
\[
    x^0=(1,p^0),\qquad z^0=(1,0).
\]
Then the following properties hold:
\begin{enumerate}[label=\textup{(\roman*)}]
\item The cluster set of the primal sequence is given by
\[
    \operatorname{cluster}(x^k)
    =
    \{1\}\times\mathbb S^1 .
\]
Moreover, every cluster point of the primal sequence is an exact stationary point of~\eqref{eq:PP} in the sense of~\eqref{eq:reduced-stationarity}.

\item The generated trajectory has infinite length, namely,
\[
    \sum_{k=0}^{\infty}\|x^{k+1}-x^k\|
    =
    +\infty .
\]
The primal sequence does not converge.

%\item Every cluster point of the generated sequence is an exact limiting
%lifted stationary point of (\ref{eq:PP}). In particular, for every
%\(\bar x\in\operatorname{cluster}(x^k)\), it satisfies the limiting lifted stationarity condition
%\eqref{eq:reduced-lifted-stationarity}.
\end{enumerate}
\end{corollary}
\begin{proof}
Since \(f\equiv1\), the S-FSPS primal and dual updates reduce to
\eqref{eq:reduced-x}--\eqref{eq:reduced-z}.
The conclusions  follow directly from Theorem~\ref{thm:main-counterexample}.
\end{proof}
%\begin{proof}The conclusions follow directly from Theorem \ref{thm:main-counterexample} and $f\equiv 1$. \end{proof}
%The exact-realization argument in the proof of
%Theorem~\ref{...} gives \(x^k=(1,p^k)\) and \(z^k=(1,0)\) for every
%\(k\). Hence, the cluster-set and infinite-length conclusions follow
%from Step~2 of that proof. For every \(x=(1,q)\), \(q\in\mathbb S^1\),
%we have
%\[
%z=(1,0)\in\partial g(Ax),\qquad
%\nabla h(x)=0,\qquad
%-(1,0,0)\in N_{\mathcal S}(x).
%\]
%Since \(A^*z=(1,0,0)\), condition~\eqref{...} holds.\end{proof}
%\begin{remark}
%The counterexample constructed in the proof of
%Theorem~\ref{thm:main-counterexample} can be specialized to the reduced
%iteration \eqref{eq:reduced-x}--\eqref{eq:reduced-z}. For the resulting
%instance, the primal sequence fails to converge, despite the existence of
%exact stationary cluster points. Therefore, the obstruction identified
%in this paper is not specific to fractional programming; it is already
%present in variable-smoothing full-splitting methods for nonconvex
%composite  optimization.
%\end{remark}


\section{Conclusion}\label{sec:conclusion}


This paper shows that the standing assumptions of S-FSPS do not
guarantee convergence of the entire primal sequence. Using a
compatible first-order jet and a $C^{1,1}$ Whitney extension theorem,
we construct two admissible instances whose primal iterates follow
slowly rotating spirals, each with cluster set
$\{1\}\times\mathbb S^1$. The resulting sequences have infinite
length and fail to converge, although every cluster point is an exact
limiting lifted stationary point.  This phenomenon persists when the feasible set is
full-dimensional, the linear operator  $A$ has full row rank,
and both the numerator and the denominator are nonconstant. The construction
also yields a nonconvergent example for the corresponding variable-smoothing,
single-loop, full-splitting method for  nonconvex composite optimization,
despite the existence of a subsequence converging to an exact stationary
point. Thus, the existing convergence theory is sharp with respect to
whole-sequence convergence: subsequences may converge to exact stationary
points while the full sequence fails to converge. The counterexamples do not arise from unboundedness or from
nonstationarity of their cluster points; rather, they show that
stationarity of all cluster points is insufficient to control the
global geometry of the generated trajectories.
%If  the set of lifted stationary points at the
%limiting objective value is finite, or  isolated, then the whole-sequence convergence follows.
%Identifying minimal conditions that ensure whole-sequence convergence while preserving the single-loop, fully splitting structure remains an important direction for future research.

%\paragraph{Why standard KL arguments do not apply directly.}



%Moreover, our construction assumes only the $C^{1,1}$ regularity required
%by the algorithmic framework. It does not establish that the resulting
%objective is semialgebraic, definable, or otherwise possesses the KL
%property. Hence the counterexample should be viewed as demonstrating that,
%without additional structural assumptions and without a compatible
%descent--relative-error mechanism, whole-sequence convergence cannot be
%expected. Finally, we emphasize that a continuum of stationary points is
%not, by itself, incompatible with the KL property. The essential obstacle
%is the lack of a KL-type finite-length argument rather than the existence
%of infinitely many stationary points.


%Nevertheless, the counterexample constructed in the proof of
%Theorem~\ref{thm:main-counterexample} can be specialized to the reduced
%iteration \eqref{eq:reduced-x}--\eqref{eq:reduced-z}. For the resulting
%instance, the primal sequence fails to converge, despite the existence of
%exact stationary cluster points. Therefore, the obstruction identified
%in this paper is not specific to fractional programming; it is already
%present in smoothing-based, single-loop, fully splitting methods for
%composite nonconvex optimization.


\section*{Statements and Declarations}

\paragraph{Funding.} The research of Min Tao was supported  by the National Natural Science Foundation of China (Grant No. 12471289).
% Insert the accurate funding statement.

\paragraph{Competing interests.}
The author has no relevant financial or non-financial interests to disclose.

\paragraph{Data availability.}
No data were generated or analyzed in this study.

\begin{thebibliography}{99}

%\bibitem{AttouchBolteSvaiter2013}
%H.~Attouch, J.~Bolte, and B.~F.~Svaiter,
%Convergence of descent methods for semi-algebraic and tame problems:
%proximal algorithms, forward--backward splitting, and regularized
%Gauss--Seidel methods,
%\emph{Math. Program.} 137 (2013), 91--129.
%
%\bibitem{BolteSabachTeboulle2014}
%J.~Bolte, S.~Sabach, and M.~Teboulle,
%Proximal alternating linearized minimization for nonconvex and nonsmooth
%problems,
%\emph{Math. Program.} 146 (2014), 459--494.
\bibitem{AB}  H. Attouch and J. Bolte, {\em On the convergence of the proximal algorithm for nonsmooth functions involving analytic features,} Math. Program., 116 (2009), pp. 5--16.

\bibitem{ABS} H.  Attouch, J. Bolte, and B.~F. Svaiter, {\em Convergence of descent methods for semi-algebraic and tame problems: proximal algorithms, forward-backward splitting, and regularized Gauss-Seidel methods}, Math. Program.,  137 (2013), pp. 91--129.
\bibitem{BeckTeboulle2012}
A.~Beck and M.~Teboulle,
Smoothing and first order methods: A unified framework,
SIAM J. Optim., 22(2) (2012), pp.~557--580.

\bibitem{BianChen2020}
W.~Bian and X.~Chen,
A smoothing proximal gradient algorithm for nonsmooth convex regression
with cardinality penalty,
SIAM J. Numer. Anal., 58(1) (2020), pp.~858--883.

\bibitem{BoltePauwels2021}
J.~Bolte and E.~Pauwels,
{\em Curiosities and counterexamples in smooth convex optimization},
Math. Program., 195 (2022), pp. 553--603.



\bibitem{BolteCombettesPauwels2024}
J.~Bolte, C.~W.~Combettes, and E.~Pauwels,
The iterates of the Frank--Wolfe algorithm may not converge,
Math. Oper. Res., 49(4) (2024), pp.~2565--2578.


\bibitem{BotBohm2020}
R.~I.~Bo\c{t} and A.~B\"ohm,
Variable smoothing for convex optimization problems using stochastic
gradients,
J. Sci. Comput., 85 (2020), Art.~33.

\bibitem{BotHendrich2015}
R.~I.~Bo\c{t} and C.~Hendrich,
A variable smoothing algorithm for solving convex optimization problems,
TOP, 23(1) (2015), pp.~124--150.

%\bibitem{BotDaoLi2022}
%R.~I. Bo\c{t}, M.~N. Dao, and G.~Li,
%{\em Extrapolated proximal subgradient algorithms for nonconvex and nonsmooth
%fractional programs},
%Math. Oper. Res., 47(3) (2022), pp. 2415--2443.

\bibitem{BotLiTao2025}
R.~I. Bo\c{t}, G.~Li, and M.~Tao,
{\em Full splitting algorithms for fractional programs with structured numerators
and denominators},
SIAM J. Optim., 35(4) (2025), pp. 2623--2653.

\bibitem{BohmWright2021}
A.~B\"ohm and S.~J.~Wright,
Variable smoothing for weakly convex composite functions,
J. Optim. Theory Appl., 188 (2021), pp.~628--649.



\bibitem{Chen2012}
X.~Chen,
{\em Smoothing methods for nonsmooth, nonconvex minimization},
Math. Program. Ser. B, 134(1) (2012), pp. 71--99.

\bibitem{DaniilidisEtAl2018}
A.~Daniilidis, M.~Haddou, E.~Le Gruyer, and O.~Ley,
{\em Explicit formulas for \(C^{1,1}\) Glaeser--Whitney extensions of 1-Taylor
fields in Hilbert spaces},
Proc. Amer. Math. Soc., 146(10) (2018), pp. 4487--4495.

%\bibitem{KumeYamada2025}
%K.~Kume and I.~Yamada,
%{\em A Variable Smoothing for Weakly Convex
%Composite Minimization with Manifold
%Constraint via Parametrization},
%available at \url{https://arxiv.org/pdf/2412.04225}.
\bibitem{KumeYamada2025}
K.~Kume and I.~Yamada,
{\em A variable smoothing for weakly convex composite minimization
with manifold constraint via parametrization},
\newblock arXiv:2412.04225v4, 2026.

\bibitem{LiuXia2024}
Y.~Liu and F.~Xia,
Proximal variable smoothing method for three-composite nonconvex nonsmooth
minimization with a linear operator,
Numer. Algorithms, 96(1) (2024), pp.~237--266.

\bibitem{LongEtAl2026}
X.-J.~Long, K.~Zeng, G.-X.~Li, M.~N.~Dao, and Z.-Y.~Peng,
Variable smoothing alternating proximal gradient algorithm for coupled
composite optimization,
J. Optim. Theory Appl., 209 (2026), Art.~19.




\bibitem{Nesterov2005}
Y.~Nesterov,
Smooth minimization of non-smooth functions,
Math. Program., 103 (2005), pp.~127--152.

\bibitem{RockWets}
 R.~T. Rockafellar and R.~J.~B. Wets, {\em Variational analysis}, Springer, Berlin, 1998.

\bibitem{SabachTeboulleVoldman2018}
S.~Sabach, M.~Teboulle, and S.~Voldman,
{\em A smoothing alternating minimization-based algorithm for clustering with
sum-min of Euclidean norms},
Pure Appl. Funct. Anal., 3(4) (2018), pp. 653--679.

\bibitem{Sun1}
 T. Sun, {\em Sequence convergence of inexact nonconvex and nonsmooth algorithms with more realistic assumptions},
Numer. Funct. Anal. Optim., 42(2)  (2021), pp. 234--250.

%\bibitem{ZhangLi2022}
%N.~Zhang and Q.~Li,
%{\em First-order algorithms for a class of fractional optimization problems},
%SIAM J. Optim., 32(1) (2022), pp. 100--129.






\end{thebibliography}



\appendix
\section{Proof of Lemma \ref{lem:uniform-moreau-estimates}}\label{appA}

\begin{proof}
Fix $w\in A(\mathcal S)$. Since $\partial g(w)$ is a nonempty closed convex
set, there exists $v_w\in\partial g(w)$ satisfying
$
\|v_w\|=\dist(0,\partial g(w))\leq\ell.$
Set $p:=\prox_{\gamma g}(w)$ and
$q:=\nabla g_\gamma(w)=(w-p)/\gamma$. Because
$q\in\partial g(p)$ and $v_w\in\partial g(w)$, monotonicity of $\partial g$
gives
$
\langle q-v_w,p-w\rangle\geq0.$
Since $p-w=-\gamma q$, it follows that
$
\|q\|^2\leq\langle v_w,q\rangle
\leq\|v_w\|\,\|q\|,$
which proves (i). (ii) follows immediately from
$p-w=-\gamma q$.
The subgradient inequality gives
$
g(u)\geq g(w)+\langle v_w,u-w\rangle.
$
Consequently,
\begin{eqnarray*}
g_\gamma(w)
\geq g(w)+\inf_{d\in\R^s}
\left\{\langle v_w,d\rangle+\frac{1}{2\gamma}\|d\|^2\right\}
=g(w)-\frac{\gamma}{2}\|v_w\|^2
\geq g(w)-\frac{\gamma\ell^2}{2}.
\end{eqnarray*}
Together with $g_\gamma(w)\leq g(w)$, this proves
(iii).
Finally, let $u,v\in A(\mathcal S)$, and choose
$v_u\in\partial g(u)$ and $v_v\in\partial g(v)$ with
$\|v_u\|,\|v_v\|\leq\ell$. Applying the subgradient inequality at $u$ and
at $v$ gives
$
g(v)\geq g(u)+\langle v_u,v-u\rangle,$ and
$
g(u)\geq g(v)+\langle v_v,u-v\rangle.
$
These two inequalities imply (iv).
\end{proof}



\section{Proof of Theorem \ref{PriTheo2R}}\label{appB}
\begin{proof}In \cite[Theorem~4.3]{BotLiTao2025}, the additional condition
$
    A(\mathcal S)\subseteq\operatorname{int}(\operatorname{dom}g)
$
is imposed only for assertions \textup{(vi)}--\textup{(vii)};
the proofs of assertions \textup{(i)}--\textup{(v)} do not use
this condition. It therefore remains  to reprove
\textup{(vi)} and \textup{(vii)}.\\
\textup{(vi)}
Let $(\overline x,\overline y,\overline z)\in\Omega$ and choose a subsequence such that
$
(x^{k_j},y^{k_j},z^{k_j})\longrightarrow(\overline x,\overline y,\overline z).
$
Set $\widehat\gamma_j:=\gamma_{k_j-1}$. Then $\widehat\gamma_j\to0$ and
$Ax^{k_j},A\overline x\in A(\mathcal S)$. By
Lemma~\ref{lem:uniform-moreau-estimates} (iii) and (iv),
\begin{align}
|g_{\widehat\gamma_j}(Ax^{k_j})-g(A\overline x)|
&\leq
|g_{\widehat\gamma_j}(Ax^{k_j})-g(Ax^{k_j})|
+|g(Ax^{k_j})-g(A\overline x)|\nonumber\\
&\leq\frac{\widehat\gamma_j\ell^2}{2}
+\ell\|Ax^{k_j}-A\overline x\|
\longrightarrow0.
\label{eq:moreau-value-convergence}
\end{align}
Since
$
\Psi(x^{k_j},z^{k_j};\widehat \gamma_j)
=g_{\widehat \gamma_j}(Ax^{k_j})+h(x^{k_j})
$
and
$
\theta_{k_j}
=\frac{\Psi(x^{k_j},z^{k_j};\widehat \gamma_j)}
{f(Kx^{k_j})},
$
letting $j\to\infty$ and using the continuity of $h$ on $\mathcal S$ and of $f$ on
$K(\mathcal S)$ yields
$
\frac{g(A\overline x)+h(\overline x)+\iota_{\mathcal S}(\overline x)}
{f(K\overline x)}=\overline\theta.
$
This proves \textup{(vi)}.

\textup{(vii)}
Let \(\{{ x}^{k_j}\}\) be a subsequence of $\{ x^k\}$ such that
\begin{eqnarray}\label{add11}
\lim\limits_{j \to +\infty} \delta_{k_j} \|{ x}^{k_j+1} - { x}^{k_j}\| = 0
 \end{eqnarray}
 and let $\overline{ x} \in {\cal S}$ be a cluster point of it. Then there exists a further subsequence $\{{ x}^{k_s}\}$ of $\{{ x}^{k_j}\}$ that converges to ${\overline{ x}}$ as $s \to +\infty$. Since $\delta_k\ge\chi L_{\nabla h}$ for all $k \geq 0$, we have $\lim_{s\to+\infty}\|{ x}^{k_s+1}-{ x}^{k_s}\|=0$, thus ${ x}^{k_s+1} \rightarrow \overline { x}$ as $s \to +\infty$.
By assertion (ii), $\{y^k\}$ is bounded. Passing to a further
subsequence if necessary, suppose that
$y^{k_s+1}\to\overline y$. Since
$y^{k_s+1}\in\partial f(Kx^{k_s})$ and
$Kx^{k_s}\to K\overline x$, the closedness of
$\operatorname{gph}\partial f$ yields
$\overline y\in\partial f(K\overline x)$.
From  the $ x$-update in (\ref{eq:x-update}) and $\nabla(g_{\gamma_{k_s-1}} \circ A)({x}^{k_s})=A^* \nabla g_{\gamma_{k_s-1}}(A {x}^{k_s})$, for all $s \geq 0$, there exists ${ \xi}^{k_s+1} \in \partial \iota_{\cal S}({ x}^{k_s+1})$ such that
\begin{eqnarray}\label{KKT9}
\ \ \ \ \ \ \ \ 0 = { \xi}^{k_s+1}  +\delta_{k_s}({ x}^{k_s+1}-{ x}^{k_s}) + A^*\nabla g_{\gamma_{k_s-1}}(A{ x}^{k_s}) + \nabla h({ x}^{k_s})  - \theta_{k_s} K^* { y}^{k_s+1}.
\end{eqnarray}

Recall
$\widehat \gamma_s:=\gamma_{k_s-1},$  and let $q_s:=\nabla g_{\widehat\gamma_s}(Ax^{k_s}),$ and $p_s:=\prox_{\widehat\gamma_s g}(Ax^{k_s}).$
By  Lemma \ref{lem:uniform-moreau-estimates}(i), $\|q_s\|\leq\ell$. Passing to a
subsequence, suppose that $q_s\to\overline q$. Moreover,
$
\|p_s-Ax^{k_s}\|
\leq\widehat\gamma_s\ell\longrightarrow0,
$
so $p_s\to A {\overline x}$. Since \(q_s\in\partial g(p_s)\), \(p_s\to A{\overline x}\), and \(q_s\to \overline{q}\),
the closedness of \(\operatorname{gph}\partial g\) yields
\(\overline q\in\partial g(A\overline{x})\).
All terms in \eqref{KKT9}, except possibly
$\xi^{k_s+1}$, are bounded; hence $\{\xi^{k_s+1}\}$ is bounded. Passing
to a further subsequence $\xi^{k_s+1}\to \overline \xi$, and using the closedness of
$\operatorname{gph}(\partial\iota_{\mathcal S})$, we obtain
$\overline\xi\in\partial\iota_{\mathcal S}(\overline x)$.
Moreover, combining with $\theta_{k_s}\to\overline \theta$ due to (iv), (\ref{add11}) and (\ref{KKT9}), we have
$
0=\overline\xi+A^*\overline q+\nabla h(\overline x)
-\overline\theta K^*\overline y.
$
Therefore,
$
0\in\partial\iota_{\mathcal S}(\overline x)
+A^*\partial g(A\overline x)+\nabla h(\overline x)
-\overline\theta K^*\partial f(K\overline x).
$
Moreover, assertion~\textup{(vi)} gives
\[
    \overline\theta=\frac{g(A\overline x)+h(\overline x)}{f(K\overline x)}.
\]
Multiplying the preceding inclusion by \(f(K\overline x)>0\) therefore shows
that \(\overline x\) is a limiting lifted stationary point of
\eqref{eq:fractional-problem}.
\end{proof}
\section{Proof of the chord estimate}\label{app:chord}
\begin{lemma}
\label{lem:chord-lower-bound}
For every \(k\ge 0\),
\[
\|p^{k+1}-p^k\|
\ge \frac{\beta_k}{\pi}.
\]
\end{lemma}
\begin{proof}
Since the angle between $u(\varphi_{k+1})$ and $u(\varphi_k)$ is $\beta_k$,
\[
\langle u(\varphi_{k+1}),u(\varphi_k)\rangle=\cos\beta_k.
\]
Consequently,
\begin{align*}
\|p^{k+1}-p^k\|^2
&=
\|r_{k+1}u(\varphi_{k+1})-r_ku(\varphi_k)\|^2\\
&=
r_{k+1}^2+r_k^2-2r_kr_{k+1}\cos\beta_k\\
&=
(r_{k+1}-r_k)^2
+4r_kr_{k+1}\sin^2\!\left(\frac{\beta_k}{2}\right).
\end{align*}
Dropping the first nonnegative term gives
\[
\|p^{k+1}-p^k\|
\ge
2\sqrt{r_kr_{k+1}}
\sin\!\left(\frac{\beta_k}{2}\right).
\]
Since $r_k,r_{k+1}\ge 1/2$,
$
2\sqrt{r_kr_{k+1}}\ge1,
$
and therefore
$
\|p^{k+1}-p^k\|
\ge
\sin\!\left(\frac{\beta_k}{2}\right).
$
For $0\le t\le\pi/2$, concavity of $\sin$ gives
$
\sin t\ge \frac{2}{\pi}t.
$
Because $0<\beta_k\le1<\pi$ for all $k\ge0$, we have
$0<\beta_k/2<\pi/2$, and hence
\[
\sin\!\left(\frac{\beta_k}{2}\right)
\ge
\frac{\beta_k}{\pi}.
\]
Thus
$
\|p^{k+1}-p^k\|\ge\frac{\beta_k}{\pi}.
$
\end{proof}

\section{Proof of the elementary spiral estimates}\label{app:spiral-estimates}

\begin{lemma}
\label{lem:elementary}
Fix \(\eta\in(0,1/2]\) and choose \(K_0\) such that
\eqref{eq:K0-conditions} holds. Then there exist absolute constants
\(\widehat C_1,\widehat C_2,\widehat C_3,\widehat C_4>0\), independent
of \(k\), \(K_0\), and \(\eta\), such that
\begin{align}
    0<d_k-d_{k+1}
    &\leq \widehat C_1\alpha_kd_k^3,
    \label{eq:d-first}\\
    \norm{v_k}
    &\leq \widehat C_2\eta d_k^2,
    \label{eq:v-size}\\
    \norm{v_{k+1}-v_k}
    &\leq \widehat C_3\eta\,\norm{p^{k+1}-p^k},
    \label{eq:v-consecutive}\\
    0\leq a_k
    &\leq \widehat C_4\eta^2d_k^2.
    \label{eq:a-size}
\end{align}
\end{lemma}
\begin{proof}
Set
$
    u_k:=u(\varphi_k),
    \;
    e_k:=(-\sin\varphi_k,\cos\varphi_k),
    \;
    \beta_k:=\varphi_{k+1}-\varphi_k
             =\eta\alpha_kd_k^2.
$
Choose \(K_0\) such that
\eqref{eq:K0-conditions} holds. Then, for every \(k\geq0\),
\[
r_k\geq\frac12,
\qquad
0\leq\beta_k\leq1,
\qquad
d_k^2\leq\eta,
\qquad
d_k^3\leq\eta^2.
\label{eq:app-K-conditions}
\]
%Indeed, \(d_k\leq d_0\leq\eta^2\) and \(0<\eta\leq1\).


\noindent To prove \eqref{eq:d-first}, define
\[
d(t):=(\log t)^{-1/2},\qquad t\geq K_0.
\]
Since
\[
d'(t)
=-\frac{1}{2t(\log t)^{3/2}}
=-\frac{d(t)^3}{2t},
\]
and \(d_k=d(k+K_0)\), the mean-value theorem implies that, for some
\(\xi_k\in(k,k+1)\),
\begin{align*}
0<d_k-d_{k+1}
&=-d'(\xi_k+K_0)\\
&=\frac{1}{2(\xi_k+K_0)
[\log(\xi_k+K_0)]^{3/2}}\\
&\leq
\frac{1}{2(k+K_0)[\log(k+K_0)]^{3/2}}\\
&\leq
\frac{2}{2(k+K_0+1)[\log(k+K_0)]^{3/2}}\\
&=2\alpha_k d_k^3,
\end{align*}
where the second inequality follows from \(k+K_0+1\leq 2(k+K_0)\).
Thus, \eqref{eq:d-first} holds with
\(\widehat C_1=2\).


Note that
\begin{equation}
    \left|d_{k+1}^2-d_k^2\right|=\left|d_{k+1}-d_k\right|\left|d_{k+1}+d_k\right|\le 2d_k\left|d_{k+1}-d_k\right|
    \leq 4\alpha_kd_k^4.
    \label{eq:app-d-square}
\end{equation}

Define the scaled radial increment
\[
    \rho_k
    :=\frac{r_{k+1}-r_k}{\alpha_k}
    =\frac{d_k-d_{k+1}}{\alpha_k}.
\]
We have that
\begin{equation}
    0\leq \rho_k\leq 2d_k^3.
    \label{eq:app-rho-size}
\end{equation}
%Indeed,
%\begin{eqnarray*}
%\rho_k
%&=&2(k+K_0+1)
%\left(
%\frac{1}{\sqrt{\log(k+K_0)}}
%-\frac{1}{\sqrt{\log(k+K_0+1)}}
%\right)\\
%&\leq&
%2(k+K_0+1)
%\frac{\log\left(1+\frac{1}{k+K_0}\right)}
%     {[\log(k+K_0)]^{3/2}}\\
%&\leq&
%2(k+K_0+1)
%\frac{1}{(k+K_0)[\log(k+K_0)]^{3/2}}\\
%&\leq&
%\frac{4}{[\log(k+K_0)]^{3/2}}
%=4d_k^3,
%\end{eqnarray*}
%where the second-to-last inequality follows from $\log(1+x)\le x$ when $x\in[0,1)$.
Define
$
    F(t):=2(t+1)\bigl(d(t)-d(t+1)\bigr).$
Then \(\rho_k=F(k+K_0)\).  Since
\[
    d(t)-d(t+1)
    =\frac12\int_t^{t+1}
      \frac{ds}{s(\log s)^{3/2}},
\]
we can write
\[
    F(t)=(t+1)\int_t^{t+1}f(s)\,ds,
    \qquad
    f(s):=\frac{1}{s(\log s)^{3/2}}.
\]
Differentiation gives
\[
    F'(t)
    =\int_t^{t+1}\bigl(f(s)+(t+1)f'(s)\bigr)\,ds,
\]
where
\[
    f'(s)
    =-\frac{1}{s^2(\log s)^{3/2}}
     -\frac{3}{2s^2(\log s)^{5/2}}.
\]
Using the preceding representation of $F'(t)$, together with the mean-value theorem, there exists $\widehat s\in(t,t+1)$ such that
\begin{eqnarray*}
|F'(t)|
&\leq& \frac{1}{\widehat s(\log(\widehat s))^{3/2}}
+(t+1)\left(\frac{1}{\widehat s^2(\log(\widehat s))^{3/2}}
+\frac{3}{2\widehat s^2(\log(\widehat s))^{5/2}}\right)\nn\\
&\leq&\frac{1}{t(\log(t))^{3/2}}
+(t+1)\left(\frac{1}{t^2(\log(t))^{3/2}}
+\frac{3}{2t^2(\log(t))^{5/2}}\right)\nn\\
&\leq&\frac{3}{t(\log(t))^{3/2}}.
\end{eqnarray*}

Since \(t\ge e^{4}\), we have
$
\frac{t+1}{t}
=1+\frac1t
\le
\frac43,
\;
1+\frac{3}{2\log t}
\le
1+\frac38
=
\frac{11}{8}.$
Hence,
$
1+\frac{t+1}{t}
\left(
1+\frac{3}{2\log t}
\right)
\le
1+\frac43\cdot\frac{11}{8}
=
\frac{17}{6}
<
3,
$
which implies
$
|F'(t)|
\le
\frac{17}{6t(\log t)^{3/2}}
<
\frac{3}{t(\log t)^{3/2}}.
$
%where the last inequality holds for $t\ge e^4$. Indeed, since $\log(t)\ge 4$, we have
%\[
%\frac{1}{t^2(\log(t))^{5/2}}
%\leq
%\frac{1}{t^2(\log(t))^{3/2}}.
%\]
%Recalling that \(\rho_k=F(k+K_0)\) and applying the preceding estimate, we obtain
%\begin{equation}
%    |\rho_{k+1}-\rho_k|
%    \leq C'_3\alpha_kd_k^3.
%    \label{eq:app-rho-difference}
%\end{equation}

By the mean-value theorem, there exists
$\zeta_k\in(k+K_0,k+K_0+1)$ such that
\begin{eqnarray*}
 &&|\rho_{k+1}-\rho_k|
 =|F(k+K_0+1)-F(k+K_0)|
 =|F'(\zeta_k)|\nn\\
 &&
 \leq
 \frac{3}{\zeta_k(\log \zeta_k)^{3/2}}\leq\frac{3}{(k+K_0)(\log (k+K_0))^{3/2}}\leq\frac{6}{(k+K_0+1)(\log (k+K_0))^{3/2}}
 \leq 12\alpha_k d_k^3.
\end{eqnarray*}
%where we use
%\[
%\frac{1}{\zeta_k}\leq\frac{1}{k+K_0},
%\qquad
%\frac{1}{k+K_0}\leq\frac{2}{k+K_0+1}.
%\]
\noindent Next, since
\[
    u_{k+1}=\cos\beta_k\,u_k+\sin\beta_k\,e_k,
\]
we have
\begin{eqnarray}
    &&v_k=\frac{r_ku_k-r_{k+1}u_{k+1}}{\alpha_k}=A_ku_k-B_ke_k,\nn\\
    &&A_k:=-\rho_k+c_k,
    \;\;
    c_k:=\frac{r_{k+1}(1-\cos\beta_k)}{\alpha_k},
    \;\;
    B_k:=\frac{r_{k+1}\sin\beta_k}{\alpha_k}.
    \label{eq:app-v-frame}
\end{eqnarray}
%where
%\begin{equation}
%
%\end{equation}
Using \(1-\cos t\leq t^2/2\), \(|\sin t|\leq|t|\), and the definition of
\(\beta_k\), we obtain
\begin{equation}
    0\leq c_k\leq \frac{1}{2}\eta^2\alpha_kd_k^4,
    \qquad
    |B_k|\leq\eta d_k^2.
    \label{eq:app-c-B-size}
\end{equation}
Combining \eqref{eq:app-rho-size}, \eqref{eq:app-v-frame},
\eqref{eq:app-c-B-size}, and \(d_k\leq\eta/2\) and $\eta\alpha_k d_k^2\le 1$, we conclude that
\[
    \norm{v_k}
    \leq |A_k|+|B_k|
    \leq 2 d_k^3+\frac{1}{2}\eta^2 \alpha_k d_k^4 +\eta d_k^2
    \leq 2 \cdot\frac{\eta}{2}\cdot d_k^2+\frac{1}{2}\eta d_k^2+\eta d_k^2\le\frac{5}{2}\eta d_k^2.
\]
This proves \eqref{eq:v-size} by setting $\widehat C_2:=5/2$.

Next, we show (\ref{eq:v-consecutive}).
First, monotonicity of \(\alpha_k\) and \(d_k\) gives
\begin{equation}
    |c_{k+1}-c_k|
    \leq c_{k+1}+c_k
    \leq \eta^2\alpha_kd_k^4.
    \label{eq:app-c-difference}
\end{equation}
 Moreover,
\eqref{eq:d-first} and \eqref{eq:app-d-square} imply
\[
    |r_{k+2}-r_{k+1}|\leq 2 \alpha_{k+1}d_{k+1}^3
\]
and
\begin{eqnarray}\label{addkey}
    \left|
       r_{k+2}d_{k+1}^2-r_{k+1}d_k^2
    \right|
    \leq 6\alpha_kd_k^4.
\end{eqnarray}
%Consequently,
%\begin{align}
%    |B_{k+1}-B_k|
%    &\leq
%    \eta\left|
%       r_{k+2}d_{k+1}^2-r_{k+1}d_k^2
%    \right|
%    +C_7\eta d_k^2(\beta_k^2+\beta_{k+1}^2)\leq C_8\eta\alpha_kd_k^4.
%    \label{eq:app-B-difference}
%\end{align}
%Here all higher-order terms are absorbed using
%\(0<\alpha_k,d_k,\eta\leq1\).


We next derive an upper bound on $|B_{k+1}-B_k|$. Recall that
\[
B_k=\eta r_{k+1}d_k^2\sinc(\beta_k),
\qquad
B_{k+1}=\eta r_{k+2}d_{k+1}^2\sinc(\beta_{k+1}),
\]
where
\[
\sinc(t):=
\begin{cases}
\dfrac{\sin t}{t}, & t\neq 0,\\[1mm]
1, & t=0.
\end{cases}
\]
Hence
\begin{align*}
B_{k+1}-B_k
={}&
\eta r_{k+2}d_{k+1}^2\sinc(\beta_{k+1})
-\eta r_{k+1}d_k^2\sinc(\beta_k)\\
={}&
\eta\bigl(r_{k+2}d_{k+1}^2-r_{k+1}d_k^2\bigr)
+\eta r_{k+2}d_{k+1}^2
   \bigl(\sinc(\beta_{k+1})-1\bigr)
-\eta r_{k+1}d_k^2
   \bigl(\sinc(\beta_k)-1\bigr).
\end{align*}
Therefore,
\begin{align*}
|B_{k+1}-B_k|
\le{}&
\eta\left|r_{k+2}d_{k+1}^2-r_{k+1}d_k^2\right|+
\eta r_{k+2}d_{k+1}^2
\left|\sinc(\beta_{k+1})-1\right|+
\eta r_{k+1}d_k^2
\left|\sinc(\beta_k)-1\right|.
\end{align*}
Since \(0\le r_k\le 1\), \(d_{k+1}\le d_k\), and, for \(|t|\le 1\),
$
|\sinc(t)-1|\le \frac{t^2}{6},
$
we obtain
\begin{align*}
&\eta r_{k+2}d_{k+1}^2
\left|\sinc(\beta_{k+1})-1\right|
+
\eta r_{k+1}d_k^2
\left|\sinc(\beta_k)-1\right|\\
&\qquad\le
\frac{\eta}{6}
\left(
d_{k+1}^2\beta_{k+1}^2+d_k^2\beta_k^2
\right)\\
&\qquad\le
\frac{\eta d_k^2}{6}
\left(
\beta_{k+1}^2+\beta_k^2
\right).
\end{align*}
Thus,
\[
|B_{k+1}-B_k|
\le
\eta\left|r_{k+2}d_{k+1}^2-r_{k+1}d_k^2\right|
+
\frac{1}{6}\eta d_k^2
\left(\beta_k^2+\beta_{k+1}^2\right).
\]

Using (\ref{addkey})
%\[
%\left|r_{k+2}d_{k+1}^2-r_{k+1}d_k^2\right|
%\le 6\alpha_kd_k^4,
%\]
together with
\[
\beta_k=\eta\alpha_kd_k^2,
\qquad
\beta_{k+1}\le\beta_k,
\]
we further have
\begin{align*}
\frac{1}{6}\eta d_k^2(\beta_k^2+\beta_{k+1}^2)
&\le
\frac{1}{3}\eta d_k^2\beta_k^2=
\frac{1}{3}\eta^3\alpha_k^2d_k^6=
\frac{1}{3}(\eta^2\alpha_kd_k^2)\,
\eta\alpha_kd_k^4\le
\frac{1}{3}\eta\alpha_kd_k^4,
\end{align*}
where the last inequality uses
\(0<\eta\alpha_kd_k^2\le1\). Consequently,
\[
|B_{k+1}-B_k|
\le
\frac{19}{3}\eta\alpha_kd_k^4.
\]



Equations \eqref{eq:app-rho-size}, \eqref{eq:app-c-B-size} and
\eqref{eq:app-c-difference}, together with the preceding estimate for $|\rho_{k+1}-\rho_k|$, yield
\begin{align}
    |A_k|
    &\leq 2 d_k^3+\frac{\eta^2}{2}\alpha_kd_k^4,
    \label{eq:app-A-size}\\
    |A_{k+1}-A_k|
    &\leq 12\alpha_kd_k^3
          +\eta^2\alpha_kd_k^4.
    \label{eq:app-A-difference}
\end{align}
The orthonormal frames \((u_k,e_k)\) and
\((u_{k+1},e_{k+1})\) differ by a rotation through \(\beta_k\); hence
\[
    \norm{u_{k+1}-u_k}=2\sin(\beta_k/2)\leq\beta_k,
    \qquad
    \norm{e_{k+1}-e_k}=2\sin(\beta_k/2)\leq\beta_k.
\]
Using the representation \eqref{eq:app-v-frame}, we therefore obtain
\begin{align}
    \norm{v_{k+1}-v_k}
    &\leq |A_{k+1}-A_k|+|B_{k+1}-B_k|\notag\\
    &+\bigl(|A_{k+1}|+|B_{k+1}|\bigr)\beta_k\notag\\
    &\leq 12 \alpha_k d_k^3+\left(2\eta^2\alpha_k+\frac{19}{3}\eta\alpha_k\right)d_k^4+2\eta\alpha_k d_k^5+\frac{1}{2}\eta^3\alpha_k^2d_k^6.
    \label{eq:app-v-difference-preliminary}
\end{align}

Indeed, dividing the right-hand side of \eqref{eq:app-v-difference-preliminary} by
$\eta^2\alpha_kd_k^2$ and using
$d_k\leq\eta^2$, $\eta\leq 1/2$, $\alpha_k\leq 1$, and
$\eta\alpha_kd_k^2\leq 1$, we obtain
\[
    12+2\eta^4+\frac{19}{3}\eta^3+2\eta^5
    +\frac12\eta^9
    <\frac{27}{2},
\]
where the last inequality follows from $\eta\leq 1/2$.
Thus
\begin{equation}
    \norm{v_{k+1}-v_k}
    \leq \frac{27}{2}\eta^2\alpha_kd_k^2.
    \label{eq:app-v-difference-final}
\end{equation}
On the other hand, Lemma~\ref{lem:chord-lower-bound} gives
\begin{equation}
    \norm{p^{k+1}-p^k}
    \geq\frac{\beta_k}{\pi}
    =\frac{\eta\alpha_kd_k^2}{\pi}.
    \label{eq:app-step-lower-bound}
\end{equation}
Combining \eqref{eq:app-v-difference-final} and
\eqref{eq:app-step-lower-bound} proves
\[
    \norm{v_{k+1}-v_k}
    \leq \frac{27}{2}\pi\eta\norm{p^{k+1}-p^k},
\]
which is \eqref{eq:v-consecutive} by setting $\widehat C_3:= \frac{27}{2}\pi$.

\medskip
Next, we prove \eqref{eq:a-size}.
Using \eqref{eq:v-size},
\begin{align*}
    a_k=\sum_{\ell=k}^{\infty}\alpha_\ell\norm{v_\ell}^2\leq (5/2)^2\eta^2
      \sum_{\ell=k}^{\infty}
      \frac{1}{2(\ell+K_0+1)\log^2(\ell+K_0)}.
\end{align*}
The integral test gives
\[
    \sum_{\ell=k}^{\infty}
      \frac{1}{(\ell+K_0+1)\log^2(\ell+K_0)}
      \leq\sum_{n=k+K_0}^{\infty}\frac{1}{n\log^2 n}
    \leq\frac{2}{\log(k+K_0)}
    =2 d_k^2.
\]
Therefore
\[
    0\leq a_k\leq \frac{25}{4} \eta^2d_k^2,
\]
which is \eqref{eq:a-size} by setting $\widehat C_4:=\frac{25}{4} $.
\end{proof}

\section{Local half-turn  estimate}\label{app:local-whitney}

\begin{proposition}\label{prop:local-half-turn}
Suppose $i>j$ and
$
0\le \phi:=\varphi_i-\varphi_j\le\pi.
$
Then,
$$
L_{j,i}:=
\sum_{\ell=j}^{i-1}\|p^{\ell+1}-p^\ell\|
\le
\left(1+\frac{\pi}{\sqrt2}\right)\|p^i-p^j\|.
$$
Consequently,
\begin{eqnarray}\label{E1:2}
\|v_i-v_j\|
\le
\widehat C_5\eta\|p^i-p^j\|,
\end{eqnarray}
where $\widehat C_5=\frac{27}{2}\pi\left(1+\frac{\pi}{\sqrt{2}}\right)$.
\end{proposition}

\begin{proof}
Since $r_k$ is increasing,
$
p^{\ell+1}-p^\ell
=
(r_{\ell+1}-r_\ell)u_{\ell+1}
+
r_\ell(u_{\ell+1}-u_\ell).
$
Moreover,
$$
\|u_{\ell+1}-u_\ell\|
=
2\sin\!\left(\frac{\varphi_{\ell+1}-\varphi_\ell}{2}\right)
\le
\varphi_{\ell+1}-\varphi_\ell.
$$
Hence
\[
\|p^{\ell+1}-p^\ell\|
\le
r_{\ell+1}-r_\ell
+
r_i(\varphi_{\ell+1}-\varphi_\ell).
\]
Summing from $\ell=j$ to $i-1$ yields
\[
L_{j,i}
\le
r_i-r_j+r_i\phi.
\]
Define
$
D:=\|p^i-p^j\|.
$
The polar-distance identity gives
$
D^2
=
(r_i-r_j)^2
+
4r_ir_j\sin^2\!\left(\frac{\phi}{2}\right).
$
Thus \begin{eqnarray}\label{add3}r_i-r_j\le D.\end{eqnarray}
 Since
$0\le\phi/2\le\pi/2$,
$
\sin\!\left(\frac{\phi}{2}\right)\ge\frac{\phi}{\pi},
$
and therefore
$
D
\ge
\frac{2}{\pi}\sqrt{r_ir_j}\,\phi.
$
Hence
\[
r_i\phi
\le
\frac{\pi}{2}
\sqrt{\frac{r_i}{r_j}}\,D
\le
\frac{\pi}{\sqrt2}D,
\]
where we used $r_i\le1$ and $r_j\ge1/2$. Consequently,
\begin{eqnarray}\label{add4}
L_{j,i}
\le
\left(1+\frac{\pi}{\sqrt2}\right)D.
\end{eqnarray}

\noindent Finally, summing \eqref{eq:v-consecutive} from $\ell=j$ to $i-1$, and using (\ref{add3}), (\ref{add4}), we have
\[
\|v_i-v_j\|
\le
\sum_{\ell=j}^{i-1}\|v_{\ell+1}-v_\ell\|
\le
\frac{27}{2}\pi\eta
\sum_{\ell=j}^{i-1}\|p^{\ell+1}-p^\ell\|
\le
\frac{27}{2}\pi(1+\frac{\pi}{\sqrt{2}})\eta\|p^i-p^j\|.
\]
Thus, (\ref{E1:2}) holds with $\widehat C_5:=\frac{27}{2}\pi(1+\frac{\pi}{\sqrt{2}})$.
\end{proof}

\begin{proposition}[Local quadratic Whitney remainder]
\label{prop:local-quadratic}
Under the assumptions of Proposition~\ref{prop:local-half-turn},
\[
\left|
a_i-a_j-\langle v_j,p^i-p^j\rangle
\right|
\le
\widehat C_6 \eta\|p^i-p^j\|^2,
\]
where $\widehat C_6:=\frac{27}{2}\pi\left(1+\frac{\pi}{\sqrt{2}}\right)^3$.
\end{proposition}

\begin{proof}
From
\[
p^{\ell+1}-p^\ell=-\alpha_\ell v_\ell
\]
and
\[
a_{\ell+1}-a_\ell=-\alpha_\ell\|v_\ell\|^2,
\]
we obtain
\[
p^i-p^j
=
-\sum_{\ell=j}^{i-1}\alpha_\ell v_\ell,
\qquad
a_i-a_j
=
-\sum_{\ell=j}^{i-1}\alpha_\ell\|v_\ell\|^2.
\]
Consequently,
\[
a_i-a_j-\langle v_j,p^i-p^j\rangle
=
\sum_{\ell=j}^{i-1}
\alpha_\ell
\langle v_j-v_\ell,v_\ell\rangle.
\]
For every $\ell\in\{j,\ldots,i\}$, the angular separation
$\varphi_\ell-\varphi_j$ also lies in $[0,\pi]$. Thus
Proposition~\ref{prop:local-half-turn} yields
\[
\|v_j-v_\ell\|
\le
\widehat C_5\eta\|p^\ell-p^j\|.
\]
Using
\[
\alpha_\ell\|v_\ell\|
=
\|p^{\ell+1}-p^\ell\|,
\]
we have
\[
\left|
a_i-a_j-\langle v_j,p^i-p^j\rangle
\right|
\le
\widehat C_5\eta
\sum_{\ell=j}^{i-1}
\|p^\ell-p^j\|
\,
\|p^{\ell+1}-p^\ell\|.
\]
Since
$
\|p^\ell-p^j\|
\le
L_{j,i},
$
it follows that
$
\left|
a_i-a_j-\langle v_j,p^i-p^j\rangle
\right|
\le
\widehat C_5\eta L_{j,i}^2.
$
Applying Proposition~\ref{prop:local-half-turn} once more proves the claim by setting $\widehat C_6:=\frac{27}{2}\pi\left(1+\frac{\pi}{\sqrt{2}}\right)^3$.
\end{proof}


\section{Different-turn separation}

\begin{lemma}
\label{lem:separation}
If \(i>j\) and \(\varphi_i-\varphi_j\geq\pi\), then
\[
    \norm{p^i-p^j}
    \geq c\,\max\{d_i,d_j\},
    \qquad
    c:=1-e^{-\pi/2}>0.
\]
\end{lemma}

\begin{proof}
For every integer \(n\geq2\), monotonicity of
\(t\mapsto1/(t\log t)\) gives
\[
    \log\log(n+1)-\log\log n
    =
    \int_n^{n+1}\frac{dt}{t\log t}
    \geq
    \frac{1}{(n+1)\log(n+1)}
    \geq
    \frac{1}{2(n+1)\log n}.
\]
Thus, setting \(n=k+K_0\) (where $K_0$ is defined in Lemma \ref{lem:elementary}) in the preceding estimate yields
\begin{align*}
    \varphi_i-\varphi_j
    &=
    \frac{\eta}{2}
    \sum_{k=j}^{i-1}
    \frac{1}{(k+K_0+1)\log(k+K_0)}\\
    &\leq
    \eta
    \log\left(\frac{\log(i+K_0)}{\log(j+K_0)}\right).
\end{align*}
If \(\varphi_i-\varphi_j\geq\pi\), then
\[
    \frac{\log(i+K_0)}{\log(j+K_0)}
    \geq e^{\pi/\eta},
    \qquad
    \frac{d_i}{d_j}
    =
    \sqrt{\frac{\log(j+K_0)}{\log(i+K_0)}}
    \leq e^{-\pi/(2\eta)}
    \leq e^{-\pi/2}.
\]
Since \(d_i\leq d_j\), the reverse triangle inequality yields
$
    \norm{p^i-p^j}
    \geq
    \left|\norm{p^i}-\norm{p^j}\right|
    =
    d_j-d_i
    \geq
    (1-e^{-\pi/2})d_j.$
\end{proof}


\end{document}
